% =====================================================================
%  Projekt dyplomowy (praca inżynierska) – WEAIiIB AGH
%  Szablon zgodny z "Zasadami dyplomowania" (Decyzja Dziekana nr 6/2023)
%  Kompilator w Overleaf: pdfLaTeX  (Menu -> Settings -> Compiler)
%  Wersja 2026-10-29: nowy układ rozdziałów (warstwa sprzętowa / programowa)
% =====================================================================

% report  -> mamy \chapter (rozdziały), potrzebne w pracy dyplomowej
% twoside -> marginesy lustrzane: nieparzyste L=3cm, parzyste P=3cm
% openany -> rozdział zaczyna się na następnej stronie (bez pustych stron)
\documentclass[a4paper,12pt,twoside,openany]{report}

% =====================================================================
%  00_preamble.tex – ustawienia wyglądu zgodne z Zasadami dyplomowania
%  WEAIiIB AGH, rozdz. XI "Założenia edytorskie"
% =====================================================================

% --- Kodowanie i język polski ---------------------------------------
\usepackage[utf8]{inputenc}
\usepackage[T1]{fontenc}
\usepackage[polish]{babel}

% --- Czcionka: Times New Roman 12 pt (wymóg) ------------------------
% newtxtext/newtxmath = Times dla tekstu i wzorów matematyczny-ch
\usepackage{newtxtext}
\usepackage{newtxmath}
\usepackage{amsmath}
% --- Marginesy (wymóg) -----------------------------------------------
% strony nieparzyste: lewy 3 cm, pozostałe 2 cm
% strony parzyste:    prawy 3 cm, pozostałe 2 cm
% przy opcji twoside: inner = margines od strony grzbietu
\usepackage[a4paper,
            inner=3cm, outer=2cm,
            top=2cm,   bottom=2cm,
            footskip=1cm]{geometry}

% --- Odstęp między wierszami: 1,25 wiersza (wymóg) -------------------
% W Wordzie "1,25 wiersza" dla Times 12 pt to ok. 17,3 pt między liniami.
% LaTeX dla 12 pt ma domyślnie 14,5 pt, więc współczynnik 1.2 daje ten sam efekt.
\usepackage{setspace}
\setstretch{1.2}

% --- Akapity ---------------------------------------------------------
\usepackage{indentfirst}          % wcięcie także w 1. akapicie (polski zwyczaj)
\setlength{\parindent}{1.25cm}
\usepackage{microtype}            % ładniejsze dzielenie i justowanie

% --- Nagłówki rozdziałów: "1. Wstęp" zamiast "Rozdział 1" -----------
\usepackage{titlesec}
\titleformat{\chapter}[hang]{\normalfont\LARGE\bfseries}{\thechapter.}{0.5em}{}
\titlespacing*{\chapter}{0pt}{0pt}{1.5em}
\titleformat{\section}{\normalfont\Large\bfseries}{\thesection.}{0.5em}{}
\titleformat{\subsection}{\normalfont\large\bfseries}{\thesubsection.}{0.5em}{}

% --- Numery stron: na dole, po zewnętrznej stronie --------------------
\usepackage{fancyhdr}
\pagestyle{fancy}
\fancyhf{}
\fancyfoot[LE,RO]{\thepage}
\renewcommand{\headrulewidth}{0pt}
\fancypagestyle{plain}{\fancyhf{}\fancyfoot[LE,RO]{\thepage}%
  \renewcommand{\headrulewidth}{0pt}}  % strony z początkiem rozdziału
\usepackage{emptypage}           % puste strony wstawiane przez \cleardoublepage – bez numeru

% --- Grafika, tabele, podpisy ----------------------------------------
\usepackage{graphicx}
\graphicspath{{img/}}             % obrazki szukane w katalogu img/
\usepackage{booktabs}             % ładne tabele (\toprule, \midrule ...)
\usepackage{float}                % opcja [H] – obrazek dokładnie tutaj
\usepackage[font=small,labelfont=bf,labelsep=period]{caption}
\usepackage{siunitx}              % jednostki: \SI{1090}{\mega\hertz}
\sisetup{output-decimal-marker={,}}
\usepackage{enumitem}             % kontrola wyglądu list
% Listy (punktowane i numerowane): znak listy na wysokości wcięcia
% akapitu (1,25 cm), tekst zaraz za nim - jednakowo w całej pracy
\setlist[itemize,enumerate]{labelindent=\parindent, leftmargin=*}
\setlist{itemsep=0pt}

% --- Kod źródłowy (Python, bash, YAML) -------------------------------
\usepackage{xcolor}
\usepackage{listings}
\lstset{
  basicstyle=\ttfamily\small\linespread{1.0}\selectfont,
  numbers=left, numberstyle=\tiny, numbersep=8pt,
  xleftmargin=2em, framexleftmargin=1.5em,
  frame=single, breaklines=true, columns=fullflexible,
  keywordstyle=\color{blue!60!black}\bfseries,
  commentstyle=\color{green!40!black}\itshape,
  stringstyle=\color{red!60!black},
  showstringspaces=false, tabsize=4,
  captionpos=b,
  % polskie znaki w kodzie (pdfLaTeX):
  literate={ą}{{\k{a}}}1 {ć}{{\'c}}1 {ę}{{\k{e}}}1 {ł}{{\l{}}}1
           {ń}{{\'n}}1 {ó}{{\'o}}1 {ś}{{\'s}}1 {ź}{{\'z}}1 {ż}{{\.z}}1
           {Ą}{{\k{A}}}1 {Ć}{{\'C}}1 {Ę}{{\k{E}}}1 {Ł}{{\L{}}}1
           {Ń}{{\'N}}1 {Ó}{{\'O}}1 {Ś}{{\'S}}1 {Ź}{{\'Z}}1 {Ż}{{\.Z}}1
}
\renewcommand{\lstlistingname}{Listing}
\renewcommand{\lstlistlistingname}{Spis listingów}

% --- Literatura i spis treści ----------------------------------------
\usepackage[nottoc]{tocbibind}    % "Literatura" pojawi się w spisie treści
\addto\captionspolish{\renewcommand{\bibname}{Literatura}}
\usepackage{url}

\usepackage{tikz}
\usetikzlibrary{arrows.meta,positioning,calc,decorations.pathreplacing,shapes.geometric,fit,backgrounds,babel}
\usepackage{pgfplots}
\pgfplotsset{compat=1.18}
\definecolor{akcent}{RGB}{0,84,159}

\newenvironment{gdzie}{%
  \par\noindent gdzie:%
  \begin{itemize}[label={},leftmargin=2em,labelsep=0pt,itemsep=0pt,parsep=0pt,topsep=2pt]%
}{\end{itemize}\par}
\widowpenalty=10000
\clubpenalty=10000
\hyphenation{Mas-sa-chu-setts}
% --- Odnośniki klikalne w PDF (ładować jako ostatni) -----------------
\usepackage[hidelinks]{hyperref}
          % pakiety, marginesy, czcionka, odstępy

% ------------------- DANE NA STRONĘ TYTUŁOWĄ ------------------------
% Tytuł MUSI być identyczny z tematem zatwierdzonym w USOS-APD!
\newcommand{\tytulPL}{Przenośny system logowania pozycji statków powietrznych w ROS 2.0 na podstawie transponderów ADS-B}
\newcommand{\tytulEN}{Portable aircraft position logging system in ROS 2.0 based on ADS-B transponders}
\newcommand{\autor}{Maksymilian Stach}
\newcommand{\kierunek}{Automatyka i Robotyka}
\newcommand{\opiekun}{dr inż. Mateusz Komorkiewicz}
\newcommand{\rok}{2027}           % rok kalendarzowy złożenia pracy

\hypersetup{pdftitle={\tytulPL}, pdfauthor={\autor}}

\begin{document}

% ---------------------- STRONA TYTUŁOWA -----------------------------
% =====================================================================
%  Strona tytułowa – wg Załącznika nr 1 "Zasad dyplomowania" WEAIiIB
%  (wzór Z1A: Projekt dyplomowy, studia I stopnia)
% =====================================================================
\newgeometry{margin=2.5cm}        % strona tytułowa: równe marginesy jak we wzorze
\begin{titlepage}
\sffamily                         % bezszeryfowa czcionka jak we wzorze
\setstretch{1.0}
\centering

\includegraphics[width=14cm]{agh_logo}\par
\vspace{2mm}
{\large\bfseries WYDZIAŁ ELEKTROTECHNIKI, AUTOMATYKI, INFORMATYKI\\
 I INŻYNIERII BIOMEDYCZNEJ\par}

\vspace{1.5cm}
{\fontsize{18}{22}\selectfont Projekt dyplomowy\par}

\vspace{1.2cm}

{\fontsize{16}{20}\selectfont\itshape \tytulPL\par}
\vspace{0.4cm}
{\fontsize{16}{20}\selectfont\itshape \tytulEN\par}

\vfill

\begin{flushleft}
\begin{tabular}{@{}p{5cm}l@{}}
Autor:             & \autor    \\
Kierunek studiów:  & \kierunek \\
Opiekun pracy:     & \opiekun  \\
\end{tabular}
\end{flushleft}

\vspace{1.5cm}
{Kraków, \rok\par}
\vspace{1cm}
\end{titlepage}
\restoregeometry

% tył strony tytułowej pusty, żeby spis treści zaczynał się na stronie nieparzystej
\cleardoublepage


% ------------------------- SPIS TREŚCI ------------------------------
\tableofcontents
\cleardoublepage

% --------------------------- ROZDZIAŁY ------------------------------
\section*{Wstęp} \addcontentsline{toc}{section}{Wstęp}
Cel pracy, teza           % 1. Wstęp: motywacja, cel, zakres, układ
%20_Theory.tex -- Rozdział 2. Podstawy teoretyczne
\tikzset{
  samolot/.pic={
    \fill[black!80] (-0.62,-0.06) -- (0.42,-0.06)
      .. controls (0.60,-0.06) and (0.62,0.06) .. (0.42,0.06) -- (-0.62,0.06) -- cycle;
    \fill[black!80] (-0.62,0.06) -- (-0.74,0.32) -- (-0.60,0.32) -- (-0.44,0.06) -- cycle;
    \fill[black!55] (0.12,-0.02) -- (-0.18,-0.30) -- (-0.06,-0.30) -- (0.26,-0.02) -- cycle;
    \fill[black!55] (-0.50,-0.01) -- (-0.66,-0.14) -- (-0.58,-0.14) -- (-0.40,-0.01) -- cycle;
  },
  satelita/.pic={
    \fill[black!75] (-0.09,-0.09) rectangle (0.09,0.09);
    \draw[black!75] (-0.13,0) -- (-0.09,0) (0.09,0) -- (0.13,0);
    \fill[akcent!60] (-0.42,-0.06) rectangle (-0.13,0.06);
    \fill[akcent!60] (0.13,-0.06) rectangle (0.42,0.06);
  },
  radar/.pic={
    \draw[black!80, line width=0.9pt] (0,0) -- (0,0.35);
    \fill[black!80] (-0.32,0.30) .. controls (-0.12,0.62) and (0.18,0.68) .. (0.36,0.62)
      -- (0.33,0.55) .. controls (0.15,0.58) and (-0.08,0.52) .. (-0.25,0.27) -- cycle;
    \draw[black!80, line width=0.9pt] (-0.3,0) -- (0.3,0);
  },
  odbiornik/.pic={
    \draw[black!80, line width=0.9pt] (0,0.25) -- (0,0.85);
    \fill[black!80] (0,0.85) circle (0.035);
    \fill[black!80] (-0.25,0) rectangle (0.25,0.25);
  },
}


% ---------------------------------------------------------------------
\chapter{Podstawy teoretyczne}
\label{chap:teoria}
% =====================================================================


W~rozdziale przedstawiono podstawy teoretyczne zagadnień wykorzystanych w~dalszej części pracy. Omówiono system ADS-B i~budowę ramek Mode~S, sposób kodowania położenia w~wiadomościach ADS-B, odbiór sygnałów za pomocą radia definiowanego programowo, wyznaczanie pozycji i~czasu przez odbiornik GNSS oraz środowisko ROS~2, w~którym udostępniane są dane. Rozdział zamyka przegląd istniejących rozwiązań.


% ---------------------------------------------------------------------
\section{Standard ADS-B i~ramki Mode~S}
% =====================================================================
\label{sec:adsb}


Początkiem dozorowania ruchu lotniczego było wykorzystanie radaru pierwotnego (ang.~\textit{primary surveillance radar}, PSR). Radar ten wysyła krótkie impulsy radiowe i~odbiera ich odbicia od statku powietrznego. Na podstawie różnicy czasu między wysłaniem impulsu a~odebraniem echa radar wyznacza odległość do celu, a~azymut z~kąta obrotu anteny. Rozwiązanie to nie dostarcza jednak informacji o~tożsamości statku powietrznego ani o~jego wysokości~\cite{sun2021}.

Brakujące informacje dostarcza radar wtórny (ang.~\textit{secondary surveillance radar}, SSR), który wysyła zapytania na częstotliwości \SI{1030}{\mega\hertz}. Transponder statku powietrznego odpowiada na nie na częstotliwości \SI{1090}{\mega\hertz}. W~pierwszych trybach pracy, oznaczonych literami A i~C, odpowiedź transpondera zawiera w~trybie~A czterocyfrowy kod identyfikacyjny w~zapisie ósemkowym (ang.~\textit{squawk code}), a~w~trybie~C wysokość barometryczną. O~ile rozwiązanie to sprawdzało się przy niewielkim natężeniu ruchu, o~tyle wraz z~rozwojem lotnictwa ujawniły się jego ograniczenia. Gdy kilka statków powietrznych znajduje się w~tym samym kierunku względem radaru, ich odpowiedzi nakładają się na siebie, co uniemożliwia poprawny odczyt. Dodatkowo liczba dostępnych kodów, wynosząca \num{4096} ($8^4$), okazała się niewystarczająca dla obecnego ruchu lotniczego~\cite{sun2021}.

Część z~tych problemów rozwiązano, wprowadzając tryb Mode~S (ang.~\textit{Mode Select}), opracowany w~latach siedemdziesiątych XX~wieku w~Lincoln Laboratory Massachusetts Institute of Technology i~wdrażany w~kontroli ruchu lotniczego od lat dziewięćdziesiątych. Najważniejszą cechą wyróżniającą ten tryb jest wysyłanie selektywnych zapytań, kierowanych do wybranego statku powietrznego, co ograniczyło problem nakładania się odpowiedzi. Dodatkowo zamiast jednego z~\num{4096} kodów każdy transponder Mode~S ma unikalny 24-bitowy adres, nazywany adresem ICAO, co pozwala rozróżnić do \num{16777216} statków powietrznych. Odpowiedzi Mode~S mają postać ramek o~długości 56 lub 112~bitów, przy czym pierwsze 5~bitów (ang.~\textit{downlink format}, DF) jednoznacznie określa strukturę i~interpretację pozostałej części ramki~\cite{sun2021}.

W~trybie Mode~S przewidziano również tzw. rozszerzony squitter (ang.~\textit{extended squitter}). Jest to długa, 112-bitowa ramka wysyłana przez transponder samoczynnie, bez uprzedniego zapytania ze strony radaru~\cite{sun2021}. Ramki tego typu przenoszą wiadomości systemu ADS-B (ang.~\textit{Automatic Dependent Surveillance--Broadcast}). W~systemie tym statek powietrzny wyznacza własną pozycję za pomocą odbiornika nawigacji satelitarnej (GNSS), a~następnie regularnie rozgłasza ją drogą radiową wraz z~prędkością i~innymi parametrami lotu. Wiadomości ADS-B nadaje ten sam transponder Mode~S i~na tej samej częstotliwości \SI{1090}{\mega\hertz}, na której wysyłane są odpowiedzi dla radaru wtórnego. Oba systemy różnią się jednak zasadą dozorowania: w~radarze wtórnym położenie wyznacza stacja naziemna, a~w~systemie ADS-B statek powietrzny. Różnicę tę przedstawiono na rysunku~\ref{rys:ssr_adsb}.
 
\begin{figure}[htbp]
  \centering
  \begin{tikzpicture}[font=\footnotesize, >=Stealth]
    \node[anchor=west] at (0,4.35) {(a) radar wtórny};
    \draw[black!50] (0,0) -- (6.6,0);
    \pic at (1.1,0) {radar};
    \pic[scale=1.1] at (5.3,3.25) {samolot};
    \draw[->, akcent, thick] (1.50,0.80) -- node[sloped, above] {zapytanie, \SI{1030}{\mega\hertz}} (4.55,3.05);
    \draw[->, black!80, thick] (4.65,2.70) -- node[sloped, below] {odpowiedź, \SI{1090}{\mega\hertz}} (1.65,0.55);
    \node[align=center, anchor=north] at (1.1,-0.08) {radar mierzy odległość\\ i~azymut};
    \node[align=center, anchor=south] at (5.3,3.55) {transponder podaje\\ kod lub wysokość};
    \begin{scope}[xshift=7.7cm]
      \node[anchor=west] at (0,4.35) {(b) ADS-B};
      \draw[black!50] (0,0) -- (7.2,0);
      \pic at (0.9,3.55) {satelita};
      \pic at (3.6,3.90) {satelita};
      \pic at (6.3,3.55) {satelita};
      \node[black!70] at (5.45,4.25) {satelity GNSS};
      \draw[->, black!50, dashed] (1.10,3.40) -- (3.15,2.55);
      \draw[->, black!50, dashed] (3.60,3.75) -- (3.60,2.70);
      \draw[->, black!50, dashed] (6.10,3.40) -- (4.05,2.55);
      \pic[scale=1.1] at (3.6,2.35) {samolot};
      \pic at (0.9,0) {odbiornik};
      \node[anchor=north] at (0.9,-0.08) {stacja bazowa};
      \pic[scale=0.8] at (6.3,1.25) {samolot};
      \node[anchor=north, align=center] at (6.3,0.95) {inny statek\\ powietrzny};
      \draw[->, akcent, thick] (3.10,1.95) -- (1.20,0.98);
      \draw[->, akcent, thick] (4.10,1.95) -- (5.65,1.35);
      \node[align=center, akcent] at (3.60,0.75) {rozgłaszanie, \SI{1090}{\mega\hertz}:\\ pozycja, prędkość,\\ identyfikacja};
    \end{scope}
  \end{tikzpicture}
  \caption{Zasada dozorowania: (a) radar wtórny -- położenie wyznacza stacja naziemna, (b) system ADS-B -- położenie wyznacza i~rozgłasza statek powietrzny; opracowanie własne na podstawie~\cite{sun2021}}
  \label{rys:ssr_adsb}
\end{figure}
W~przestrzeni powietrznej Unii Europejskiej nadajnik ADS-B pracujący na częstotliwości \SI{1090}{\mega\hertz} w~trybie rozszerzonego squittera stanowi obowiązkowe wyposażenie statków powietrznych o~maksymalnej masie startowej powyżej \SI{5700}{kg} lub o~maksymalnej prędkości przelotowej przekraczającej \SI{250}{kt}~\cite{ue2023_1770}. W~efekcie wiadomości ADS-B nadaje obecnie większość statków powietrznych~\cite{sun2021}. Rozważano również inne sposoby realizacji ADS-B, m.in. łącze UAT (ang.~\textit{Universal Access Transceiver}). Najszerzej przyjęła się jednak realizacja oparta na rozszerzonym squitterze Mode~S, a~łącze UAT wdrożono jedynie częściowo, np. w~Stanach Zjednoczonych~\cite{sun2021}. W~dalszej części pracy pojęcie ADS-B odnosi się wyłącznie do realizacji opartej na rozszerzonym squitterze.

Ramka ADS-B ma długość 112~bitów i~jest podzielona na pięć pól, które przedstawiono na rysunku~\ref{rys:ramka_df17}~\cite{sun2021}.
 
\begin{figure}[htbp]
  \centering
  \begin{tikzpicture}[font=\footnotesize, x=0.9cm]
    \def\h{0.8}
    \foreach \a/\b/\n/\r in {0/1.1/DF/1--5, 1.1/2.1/CA/6--8, 2.1/5.6/ICAO/9--32,
                             5.6/6.8/TC/33--37, 6.8/11.6/ME (dane)/38--88, 11.6/15/PI/89--112} {
      \draw[thick, fill=akcent!8] (\a,0) rectangle (\b,\h);
      \node at ({(\a+\b)/2},{\h/2}) {\textbf{\n}};
      \node[above] at ({(\a+\b)/2},\h) {\r};
    }
    \fill[akcent!20] (5.6,0) rectangle (6.8,\h);
    \node at (6.2,{\h/2}) {\textbf{TC}};
    \draw[thick] (5.6,0) rectangle (11.6,\h);
    \draw[decorate, decoration={brace, mirror, amplitude=4pt}, black!70]
      (5.62,-0.08) -- (11.58,-0.08) node[midway, below=4pt] {pole ME (56 bitów)};
    \node[anchor=east, black!70] at (-0.15,{\h+0.22}) {bity:};
    \node[anchor=east, black!70] at (-0.15,-1.0) {hex:};
    \node[anchor=east, black!70] at (-0.15,-1.75) {znaczenie:};
    \node at (1.05,-1.0) {\texttt{8D}};
    \node at (3.85,-1.0) {\texttt{461E16}};
    \node at (8.6,-1.0)  {\texttt{231893B61E0820}};
    \node at (13.3,-1.0) {\texttt{3120B1}};
    \node[align=center] at (1.05,-1.75) {DF\,=\,17\\ CA\,=\,5};
    \node[align=center] at (3.85,-1.75) {adres ICAO\\ nadawcy};
    \node[align=center] at (8.6,-1.75)  {TC\,=\,4: identyfikacja,\\ znak wywoławczy \texttt{FIN6G}};
    \node[align=center] at (13.3,-1.75) {suma kontrolna\\ poprawna};
  \end{tikzpicture}
  \caption{Pola ramki ADS-B na przykładzie ramki \texttt{8D461E16231893B61E08203120B1} zarejestrowanej przez stację bazową; pierwsze 5~bitów pola ME stanowi kod typu TC; szerokości pól nie są proporcjonalne do liczby bitów; opracowanie własne na podstawie~\cite{sun2021}}
  \label{rys:ramka_df17}
\end{figure}
 
\newpage
Znaczenie pól ramki i przykładowe wartości zestawiono w~tabeli~\ref{tab:pola_ramki}.
\begin{table}[H]
  \centering
  \caption{Pola ramki ADS-B i interpretacja ich wartości; opracowanie własne na podstawie~\cite{sun2021}}
  \label{tab:pola_ramki}
  \small
  \begin{tabular}{@{}lcp{10.2cm}@{}}
    \toprule
    Pole & Bity & Znaczenie i przykładowe wartości \\
    \midrule
    DF & 1--5 & format ramki: 17 -- wiadomość ADS-B nadana przez transponder Mode~S, 18 -- wiadomość ADS-B z~urządzenia niebędącego transponderem, 11 -- krótki squitter, 4, 5, 20, 21 -- odpowiedzi na zapytania radaru wtórnego \\
    CA & 6--8 & możliwości transpondera: 0 -- poziom~1 (tylko dozorowanie), 4 -- poziom~2 lub wyższy, statek powietrzny na ziemi, 5 -- poziom~2 lub wyższy, statek powietrzny w~powietrzu, 6 -- poziom~2 lub wyższy, na ziemi lub w~powietrzu \\
    ICAO & 9--32 & 24-bitowy adres statku powietrznego; adresy są przydzielane państwom blokami, dlatego początek adresu wskazuje państwo rejestracji \\
    ME & 33--88 & 56 bitów danych; pierwsze 5~bitów to kod typu TC, który określa, jak interpretować pozostałe 51~bitów (położenie, prędkość, identyfikacja i~inne) \\
    PI & 89--112 & 24-bitowa suma kontrolna CRC obliczona z~pierwszych 88~bitów ramki \\
    \bottomrule
  \end{tabular}
\end{table}
W~przykładzie z~rysunku~\ref{rys:ramka_df17} wartość DF\,=\,17 oznacza wiadomość ADS-B. Wartość CA\,=\,5 wskazuje na transponder poziomu~2 lub wyższego, a~ponadto informuje, że statek powietrzny znajduje się w~powietrzu~\cite{sun2021}.

Pole ME ramki z~rysunku~\ref{rys:ramka_df17} zapisano szesnastkowo jako \texttt{231893B61E0820}. Pierwszy bajt, \texttt{23}, ma postać binarną \texttt{00100011}. Jego pierwsze pięć bitów (\texttt{00100}) daje kod typu TC\,=\,4, czyli wiadomość identyfikacyjną. Kolejne trzy bity (\texttt{011}) zawierają kategorię statku powietrznego. Według~\cite{sun2021} wartość 3 dla TC\,=\,4 oznacza średni statek powietrzny kategorii~2, czyli o~masie od 34 do 136~t. Pozostałe 48~bitów dzieli się na osiem bloków po 6~bitów, z~których każdy koduje jeden znak znaku wywoławczego. Wartości liczbowe bloków odpowiadają znakom zgodnie z~tablicą podaną  w~książce~\cite{sun2021}: liczby od 1 do 26 oznaczają litery A--Z, liczba 32 -- spację, a~liczby od 48 do 57 -- cyfry 0--9. Przebieg dekodowania przedstawiono w~tabeli~\ref{tab:znak_wywolawczy}. Wynikiem jest znak wywoławczy FIN6G, dopełniony trzema spacjami do ośmiu znaków.

\begin{table}[htbp]
  \centering
  \caption{Dekodowanie znaku wywoławczego z~pola ME ramki z~rysunku~\ref{rys:ramka_df17}; symbol \textvisiblespace{} oznacza spację; opracowanie własne na podstawie~\cite{sun2021}}
  \label{tab:znak_wywolawczy}
  \small
  \begin{tabular}{@{}lcccccccc@{}}
    \toprule
    Numer bloku & 1 & 2 & 3 & 4 & 5 & 6 & 7 & 8 \\
    \midrule
    Bity    & \texttt{000110} & \texttt{001001} & \texttt{001110} & \texttt{110110} & \texttt{000111} & \texttt{100000} & \texttt{100000} & \texttt{100000} \\
    Wartość & 6 & 9 & 14 & 54 & 7 & 32 & 32 & 32 \\
    Znak    & F & I & N & 6 & G & \textvisiblespace & \textvisiblespace & \textvisiblespace \\
    \bottomrule
  \end{tabular}
\end{table}
\newpage
Ostatnim elementem ramki jest 24-bitowe pole PI, zawierające cykliczną sumę kontrolną CRC (ang.~\textit{cyclic redundancy check}), obliczaną na podstawie pierwszych 88~bitów ramki~\cite{dump1090}. Odbiornik wyznacza własną wartość sumy dla odebranych danych i~porównuje ją z~wartością zapisaną w~polu PI. Niezgodność wykazana podczas weryfikacji świadczy o~uszkodzeniu ramki w~trakcie transmisji radiowej. Wykorzystany algorytm CRC umożliwia nie tylko wykrywanie przekłamań, ale również korekcję pojedynczych błędów bitowych~\cite{sun2021}.

Przy wykorzystaniu wiadomości ADS-B jako danych referencyjnych należy uwzględnić ograniczenia systemu. Położenie podaje sam statek powietrzny, a~protokół nie przewiduje uwierzytelnienia nadawcy ani szyfrowania wiadomości. System jest przez to podatny na celowe zakłócanie oraz nadawanie fałszywych wiadomości (ang.~\textit{spoofing})~\cite{strohmeier2015security}. Kanał \SI{1090}{\mega\hertz} jest współdzielony przez transpondery wszystkich statków powietrznych, które nadają zarówno wiadomości ADS-B, jak i~odpowiedzi na zapytania radarów wtórnych. Przy dużym natężeniu ruchu część ramek nakłada się więc na siebie i~zostaje utracona~\cite{strohmeier2014realities}. Adres ICAO może ponadto zostać przeprogramowany, co obserwowano m.in. w~przypadku statków powietrznych wojskowych~\cite{sun2021}. Wymagania wynikające z~tych ograniczeń sformułowano w~podrozdziale~\ref{sec:wymagania_referencja}.


% ---------------------------------------------------------------------
\section{Kodowanie położenia CPR}\label{sec:cpr}
% =====================================================================


W~wiadomości o~położeniu na każdą ze współrzędnych geograficznych przeznaczono 17~bitów~\cite{sun2021}. Bezpośredni zapis długości geograficznej z~zakresu \ang{360} na 17~bitach dawałby rozdzielczość około \ang{0.0027}, czyli na równiku około \SI{300}{m}. Rozdzielczość ta jest wielokrotnie gorsza od dokładności odbiorników GNSS. Z~tego powodu stosuje się kodowanie CPR (ang.~\textit{compact position reporting}), w~którym przesyła się położenie statku powietrznego wewnątrz strefy o~ograniczonych wymiarach. Dla strefy o~rozciągłości \ang{6} szerokości geograficznej ta sama liczba bitów zapewnia rozdzielczość około \SI{5}{m}. Ceną za dokładność jest niejednoznaczność, ponieważ pojedyncza wiadomość nie wskazuje, w~której strefie znajduje się statek powietrzny~\cite{sun2021}.

Statek powietrzny nadaje na przemian dwa rodzaje wiadomości o~położeniu: parzyste i~nieparzyste. Różnią się one rozmiarem stref szerokości geograficznej, który wynosi~\cite{sun2021}
\begin{equation}
  \Delta\varphi_i = \frac{\ang{360}}{4N_{\mathrm{Z}} - i},
  \label{eq:dlat}
\end{equation}
\begin{gdzie}
  \item $\Delta\varphi_i$ -- rozciągłość strefy szerokości geograficznej, w~stopniach,
  \item $N_{\mathrm{Z}} = 15$ -- liczba stref szerokości geograficznej między równikiem a~biegunem,
  \item $i$ -- rodzaj wiadomości: $i = 0$ dla wiadomości parzystej, $i = 1$ dla nieparzystej.
\end{gdzie}
Strefy wiadomości parzystych mają więc rozciągłość \ang{6}, a~nieparzystych -- około \ang{6.1}. Liczba stref długości geograficznej, oznaczana $\operatorname{NL}(\varphi)$, zależy od szerokości geograficznej $\varphi$. Wynosi ona 59 przy równiku i~maleje do 1 w~pobliżu biegunów~\cite{sun2021}.
Wyraża ją wzór~\cite{sun2021}
\begin{equation}
  \operatorname{NL}(\varphi) = \left\lfloor
    \frac{2\pi}{\arccos\left(1 - \dfrac{1 - \cos\dfrac{\pi}{2N_{\mathrm{Z}}}}{\cos^{2}\varphi}\right)}
  \right\rfloor,
  \label{eq:nl}
\end{equation}
w~którym symbol $\lfloor\cdot\rfloor$ oznacza zaokrąglenie w~dół do liczby całkowitej. Przebieg funkcji~\eqref{eq:nl} przedstawiono na rysunku~\ref{rys:nl}. Dla szerokości geograficznej Krakowa $\operatorname{NL} = 38$, co daje strefy o~rozciągłości około \ang{9.5} długości geograficznej.

\begin{figure}[htbp]
  \centering
  \begin{tikzpicture}
    \begin{axis}[
      width=13cm, height=6.2cm,
      xmin=0, xmax=90, ymin=0, ymax=62,
      xtick={0,10,...,90}, ytick={0,10,...,60},
      xlabel={$\varphi$ [\si{\degree}]}, ylabel={$\operatorname{NL}(\varphi)$ [--]},
      grid=major, grid style={black!12},
      tick label style={font=\footnotesize}, label style={font=\footnotesize},
      /pgf/number format/use comma,
    ]
      \addplot[const plot mark left, thick, akcent] coordinates {
        (0,59) (10.4705,58) (14.8282,57) (18.1863,56) (21.0294,55) (23.545,54) (25.8292,53)
        (27.939,52) (29.9114,51) (31.7721,50) (33.5399,49) (35.229,48) (36.8503,47) (38.4124,46)
        (39.9226,45) (41.3865,44) (42.8091,43) (44.1945,42) (45.5463,41) (46.8673,40) (48.1604,39)
        (49.4278,38) (50.6715,37) (51.8934,36) (53.0952,35) (54.2782,34) (55.4438,33) (56.5932,32)
        (57.7275,31) (58.8476,30) (59.9546,29) (61.0492,28) (62.1322,27) (63.2043,26) (64.2662,25)
        (65.3185,24) (66.3617,23) (67.3965,22) (68.4232,21) (69.4424,20) (70.4545,19) (71.4599,18)
        (72.4588,17) (73.4518,16) (74.4389,15) (75.4206,14) (76.3968,13) (77.3679,12) (78.3337,11)
        (79.2943,10) (80.2492,9) (81.198,8) (82.1396,7) (83.072,6) (83.9917,5) (84.8917,4)
        (85.7554,3) (86.5354,2) (87,1) (90,1)
      };
      \draw[dashed, black!70] (axis cs:50.06,0) -- (axis cs:50.06,38);
      \node[anchor=south west, font=\footnotesize] at (axis cs:51,39) {Kraków: $\varphi \approx \ang{50.1}$, $\operatorname{NL} = 38$};
    \end{axis}
  \end{tikzpicture}
  \caption{Liczba stref długości geograficznej w~funkcji szerokości geograficznej; opracowanie własne na podstawie wzoru~\eqref{eq:nl}}
  \label{rys:nl}
\end{figure}

W~wiadomości przesyła się położenie statku powietrznego wewnątrz strefy, wyrażone ułamkiem jej rozmiaru~\cite{sun2021}:
\begin{equation}
  \varphi_{\mathrm{cpr},i} = \frac{Y_i}{2^{17}},
  \label{eq:cpr}
\end{equation}
\begin{gdzie}
  \item $\varphi_{\mathrm{cpr},i}$ -- położenie wewnątrz strefy szerokości geograficznej, z~przedziału $[0, 1)$,
  \item $Y_i$ -- 17-bitowa liczba całkowita odczytana z~wiadomości.
\end{gdzie}
Długość geograficzną koduje się analogicznie. Pojedyncza wiadomość wskazuje więc tylko zbiór możliwych położeń, po jednym w~każdej strefie. Do wyboru właściwego z~nich stosuje się dekodowanie globalne albo lokalne, których ideę przedstawiono w~uproszczonej, jednowymiarowej postaci na rysunku~\ref{rys:cpr_idea}.

\begin{figure}[htbp]
  \centering
  \begin{tikzpicture}[font=\footnotesize, x=0.38cm, y=1cm, >=Stealth]
    \node[anchor=west] at (-6.0,3.55) {(a) dekodowanie globalne};
    \node[anchor=east] at (-0.4,2.8) {wiadomość parzysta};
    \draw[black!60] (0,2.8) -- (30,2.8);
    \foreach \x in {0,6,...,30} \draw[black!60] (\x,2.65) -- (\x,2.95);
    \foreach \x in {5,11,17,23,29} \fill[akcent] (\x,2.8) circle (2.3pt);
    \node[anchor=east] at (-0.4,2.0) {wiadomość nieparzysta};
    \draw[black!60] (0,2.0) -- (30,2.0);
    \foreach \x in {0,5,...,30} \draw[black!60] (\x,1.85) -- (\x,2.15);
    \foreach \x in {2,7,12,17,22,27} \fill[black!70] (\x,2.0) circle (2.3pt);
    \draw[dashed, thick] (17,1.7) -- (17,3.2);
    \node[anchor=south] at (17,3.2) {jedyny wspólny kandydat};
    \node[anchor=west] at (-6.0,1.15) {(b) dekodowanie lokalne};
    \fill[akcent!12] (12,0.2) rectangle (18,0.8);
    \node[anchor=east] at (-0.4,0.5) {wiadomość parzysta};
    \draw[black!60] (0,0.5) -- (30,0.5);
    \foreach \x in {0,6,...,30} \draw[black!60] (\x,0.35) -- (\x,0.65);
    \foreach \x in {5,11,23,29} \fill[akcent!35] (\x,0.5) circle (2.3pt);
    \fill[akcent] (17,0.5) circle (2.8pt);
    \draw[->, thick] (15,-0.25) -- (15,0.33);
    \node[anchor=north] at (15,-0.25) {punkt odniesienia};
    \node[anchor=south, black!70] at (15,0.8) {$\pm$ pół strefy};
    \draw[->, black!70] (0,-1.05) -- (31.2,-1.05);
    \node[anchor=north] at (15,-1.5) {położenie [j.u.]};
    \foreach \x in {0,5,...,30} {
      \draw[black!70] (\x,-1.05) -- (\x,-1.12);
      \node[anchor=north] at (\x,-1.12) {\x};
    }
  \end{tikzpicture}
  \caption{Idea dekodowania położenia CPR w~jednym wymiarze, dla stref o~rozmiarach 6 i~5~jednostek umownych (j.u.): (a) dekodowanie globalne -- położenie wyznacza jedyny wspólny kandydat wiadomości parzystej i~nieparzystej, (b) dekodowanie lokalne -- położenie wyznacza kandydat leżący bliżej niż pół strefy od punktu odniesienia; kreski oznaczają granice stref, a~kropki -- możliwe położenia; opracowanie własne na podstawie~\cite{sun2021}}
  \label{rys:cpr_idea}
\end{figure}

W~dekodowaniu globalnym wykorzystuje się parę wiadomości: parzystą i~nieparzystą. Strefy obu rodzajów wiadomości mają różne rozmiary, dlatego tylko jedno z~możliwych położeń wskazywanych przez wiadomość parzystą pokrywa się z~położeniem wskazywanym przez wiadomość nieparzystą (rysunek~\ref{rys:cpr_idea}a). Obie wiadomości powinny zostać odebrane w~krótkim odstępie czasu, w~praktyce nie dłuższym niż około \SI{10}{s}. Statek powietrzny nie może przy tym przekroczyć granicy stref, dla której zmienia się wartość $\operatorname{NL}$~\cite{sun2021}. W~dekodowaniu lokalnym wystarcza jedna wiadomość, jeżeli znane jest położenie odniesienia leżące bliżej niż pół strefy od rzeczywistego położenia (rysunek~\ref{rys:cpr_idea}b). Dla wiadomości nadawanych w~powietrzu warunek ten oznacza odległość mniejszą niż \SI{180}{NM}~\cite{sun2021}, co odpowiada połowie strefy szerokości geograficznej: $\ang{3} \cdot \SI[per-mode=symbol]{111.2}{\kilo\metre\per\degree} \approx \SI{333}{km}$. Punktem odniesienia jest zwykle wcześniej wyznaczone położenie statku powietrznego~\cite{sun2021}. W~przypadku odbiornika naziemnego może nim być również położenie odbiornika, o~ile statek powietrzny znajduje się w~odległości mniejszej niż \SI{180}{NM}.

W~dekodowaniu globalnym oblicza się najpierw indeks strefy szerokości geograficznej, a~następnie szerokość geograficzną dla każdej z~wiadomości~\cite{sun2021}:
\begin{equation}
  j = \left\lfloor 59\,\varphi_{\mathrm{cpr},0} - 60\,\varphi_{\mathrm{cpr},1} + \frac{1}{2} \right\rfloor,
  \qquad
  \varphi_i = \Delta\varphi_i \bigl( \operatorname{mod}(j,\, 60 - i) + \varphi_{\mathrm{cpr},i} \bigr),
  \label{eq:globalne}
\end{equation}
\begin{gdzie}
  \item $j$ -- indeks strefy szerokości geograficznej,
  \item $\varphi_i$ -- szerokość geograficzna wyznaczona z~wiadomości rodzaju $i$, w~stopniach,
  \item $\operatorname{mod}(a, b)$ -- reszta z~dzielenia $a$ przez $b$, nieujemna dla $b > 0$.
\end{gdzie}
Wartości nie mniejsze niż \ang{270} odpowiadają półkuli południowej i~pomniejsza się je o~\ang{360}. Jako wynik przyjmuje się szerokość geograficzną obliczoną z~nowszej wiadomości pary~\cite{sun2021}. W~dekodowaniu lokalnym indeks strefy wyznacza się z~położenia odniesienia $\varphi_{\mathrm{ref}}$~\cite{sun2021}:
\begin{equation}
  j = \left\lfloor \frac{\varphi_{\mathrm{ref}}}{\Delta\varphi_i} \right\rfloor
    + \left\lfloor \frac{\operatorname{mod}(\varphi_{\mathrm{ref}},\, \Delta\varphi_i)}{\Delta\varphi_i}
    - \varphi_{\mathrm{cpr},i} + \frac{1}{2} \right\rfloor,
  \qquad
  \varphi = \Delta\varphi_i \left( j + \varphi_{\mathrm{cpr},i} \right).
  \label{eq:lokalne}
\end{equation}
Długość geograficzną oblicza się w~obu metodach analogicznie, przy czym rozmiar strefy wynosi $\ang{360} / \max\bigl(\operatorname{NL}(\varphi) - i,\, 1\bigr)$ i~zależy od wyznaczonej wcześniej szerokości geograficznej~\cite{sun2021}.

Przykładowo niech stacja bazowa znajduje się na szerokości geograficznej $\varphi_{\mathrm{ref}} = \ang{50.0}$, a~statek powietrzny na szerokości \ang{50.3169}, jak w~przykładzie z~podrozdziału~\ref{sec:model}. Wiadomość parzysta zawiera wówczas wartość $Y_0 = \num{50613}$, czyli $\varphi_{\mathrm{cpr},0} \approx \num{0.38615}$. Ze wzoru~\eqref{eq:lokalne} otrzymuje się $j = 8 + \lfloor \num{0.33333} - \num{0.38615} + \num{0.5} \rfloor = 8$ oraz $\varphi = \ang{6} \cdot (8 + \num{0.38615}) \approx \ang{50.31688}$. Wynik różni się od rzeczywistej szerokości geograficznej o~około \SI{2}{m}, co wynika z~rozdzielczości kodowania.

Algorytm CPR nie gwarantuje poprawności każdego wyniku i~w~rzadkich przypadkach prowadzi do błędnego położenia. Z~tego powodu zaleca się sprawdzanie wiarygodności wyznaczonych położeń, np. przez porównanie ich z~zasięgiem odbiornika lub z~położeniami obliczonymi z~innych par wiadomości~\cite{sun2021}.


%---------------------------------------------------------------------------
\section{Odbiór sygnałów z~wykorzystaniem SDR}
% =====================================================================


\label{sec:sdr}

Odpowiedzi Mode~S, w~tym ramki ADS-B, są przesyłane z~szybkością \SI{1}{Mbit/s}~\cite{sun2021}. Każda ramka rozpoczyna się preambułą o~czasie trwania \SI{8}{\micro\second}, po której następuje blok danych o~długości 56 lub \SI{112}{\micro\second}. Dane zakodowano za pomocą modulacji położenia impulsu (ang.~\textit{pulse position modulation}, PPM). Bit o~wartości 1 odpowiada impulsowi w~pierwszej połowie mikrosekundy, a~bit o~wartości 0 impulsowi w~drugiej połowie~\cite{sun2021}. Na rysunku~\ref{rys:ppm} przedstawiono początek ramki z~rysunku~\ref{rys:ramka_df17}.

\begin{figure}[htbp]
  \centering
  \begin{tikzpicture}[font=\footnotesize, x=0.8cm, y=1cm]
    \draw[black!60] (0,0) -- (16.3,0);
    \foreach \t in {0,...,16} \draw[black!60] (\t,0) -- (\t,-0.08);
    \foreach \t in {0,2,...,16} \node[anchor=north] at (\t,-0.08) {\t};
    \node[anchor=north west] at (16.35,-0.08) {$t$ [\si{\micro\second}]};
    \foreach \t in {0,1,3.5,4.5} \fill[akcent] (\t,0) rectangle ++(0.5,0.8);
    \foreach \k/\b in {0/1,1/0,2/0,3/0,4/1,5/1,6/0,7/1} {
      \draw[black!25, dashed] ({8+\k},0) -- ({8+\k},1.05);
      \ifnum\b=1
        \fill[black!75] ({8+\k},0) rectangle ++(0.5,0.8);
      \else
        \fill[black!75] ({8.5+\k},0) rectangle ++(0.5,0.8);
      \fi
      \node at ({8.5+\k},1.25) {\b};
    }
    \draw[black!25, dashed] (16,0) -- (16,1.05);
    \draw[decorate, decoration={brace, amplitude=5pt}] (0,1.55) -- (8,1.55)
      node[midway, above=5pt] {preambuła (\SI{8}{\micro\second})};
    \draw[decorate, decoration={brace, amplitude=5pt}] (8,1.55) -- (16,1.55)
      node[midway, above=5pt] {pierwsze 8~bitów danych: \texttt{8D} = \texttt{10001101}};
  \end{tikzpicture}
  \caption{Preambuła i~pierwsze bity ramki z~rysunku~\ref{rys:ramka_df17} zakodowane metodą PPM; opracowanie własne na podstawie~\cite{sun2021}}
  \label{rys:ppm}
\end{figure}

W~radiu definiowanym programowo (ang.~\textit{software-defined radio}, SDR) większość operacji przetwarzania sygnału, w~tym demodulację, wykonuje oprogramowanie. Część sprzętowa ogranicza się do odbioru sygnału, przeniesienia go na niższą częstotliwość i~próbkowania. Użytkownik wybiera programowo m.in. częstotliwość środkową odbiornika i~częstotliwość próbkowania~\cite{sun2021}. Najpopularniejszymi tanimi odbiornikami SDR są urządzenia RTL-SDR, zbudowane na układzie Realtek RTL2832U i~pierwotnie przeznaczone do odbioru telewizji cyfrowej DVB-T~\cite{sun2021,rtlsdrv3}. Ich zakres pracy obejmuje częstotliwości od \SI{24}{\mega\hertz} do \SI{1766}{\mega\hertz}. Maksymalna częstotliwość próbkowania wynosi \SI{3.2}{MS/s}, przy czym stabilną pracę uzyskuje się do około \SI{2.8}{MS/s}~\cite{sun2021}. Do odbioru ramek Mode~S wymagana jest częstotliwość próbkowania co najmniej \SI{2}{MS/s}, ponieważ impuls PPM trwa \SI{0.5}{\micro\second}~\cite{sun2021}. Odbiornik RTL-SDR Blog~V3 zawiera tuner R820T2 (oznaczany także jako R860) oraz generator kwarcowy z~kompensacją temperaturową o~stabilności \SI{1}{ppm}~\cite{rtlsdrv3www}. Próbki sygnału przekazuje do komputera przez interfejs USB w~postaci par wartości I/Q o~rozdzielczości 8~bitów~\cite{dump1090}.

Dostępne są także odbiorniki SDR o~szerszym zakresie częstotliwości, wyższej częstotliwości próbkowania i~z~możliwością nadawania, np. HackRF One i~LimeSDR~\cite{sun2021}. Wybór odbiornika dla stacji bazowej uzasadniono w~rozdziale~\ref{chap:sprzet}.

Demodulację i~dekodowanie ramek wykonuje oprogramowanie. Najbardziej znanym otwartym dekoderem Mode~S jest program dump1090, napisany w~języku C przez Salvatore Sanfilippo i~udostępniony w~2013~roku na licencji BSD~\cite{sun2021}. Program pobiera z~odbiornika próbki I/Q z~częstotliwością \SI{2}{MS/s}, oblicza ich moduł, wyszukuje preambuły ramek, demoduluje bity PPM i~sprawdza sumę kontrolną, korygując pojedyncze błędy bitowe~\cite{dump1090}. W~trybie sieciowym dump1090 udostępnia przez protokół TCP ramki zapisane szesnastkowo (port~30002) oraz zdekodowane wiadomości w~tekstowym formacie SBS, nazywanym też formatem BaseStation (port~30003)~\cite{dump1090}. Program ma wiele odmian rozwijanych niezależnie, m.in. przez firmę FlightAware~\cite{sun2021}. Przebieg odbioru sygnału przedstawiono na rysunku~\ref{rys:tor_sdr}. Dekodowanie wiadomości ADS-B, w~tym położeń zakodowanych metodą CPR, umożliwia również biblioteka pyModeS dla języka Python~\cite{sun2020pymodes}.


\begin{figure}[htbp]
  \centering
  \begin{tikzpicture}[font=\footnotesize, >=Stealth,
      blok/.style={draw, thick, fill=akcent!8, align=center, minimum height=1.1cm, text width=2.0cm, inner sep=3pt},
      wyj/.style={draw, thick, fill=black!5, align=center, minimum height=0.9cm, text width=2.3cm, inner sep=3pt}]
    \node[blok, text width=2.7cm] (ant) {antena};
    \node[blok, text width=2.7cm, right=0.65cm of ant] (tun) {tuner R820T2:\\ wzmocnienie i~przemiana\\ częstotliwości};
    \node[blok, text width=2.7cm, right=0.65cm of tun] (adc) {RTL2832U:\\ przetwornik A/C,\\ próbki I/Q};
    \draw[->, thick] (ant) -- (tun);
    \draw[->, thick] (tun) -- (adc);
    \node[blok, below=1.8cm of ant.south west, anchor=north west] (mod) {obliczenie\\ modułu próbek};
    \node[blok, right=0.65cm of mod] (pre) {wykrycie\\ preambuły};
    \node[blok, right=0.65cm of pre] (ppm) {demodulacja\\ PPM};
    \node[blok, right=0.65cm of ppm] (crc) {kontrola CRC\\ i~korekcja\\ błędów};
    \node[wyj, right=0.65cm of crc, yshift=0.7cm] (avr) {port 30002:\\ ramki (hex)};
    \node[wyj, right=0.65cm of crc, yshift=-0.7cm] (sbs) {port 30003:\\ wiadomości SBS};
    \draw[->, thick] (mod) -- (pre);
    \draw[->, thick] (pre) -- (ppm);
    \draw[->, thick] (ppm) -- (crc);
    \draw[->, thick] (crc.east) -- (avr.west);
    \draw[->, thick] (crc.east) -- (sbs.west);
    \draw[->, thick] (adc.south) -- node[right] {USB, \SI{2}{MS/s}} ++(0,-0.95) -| (mod.north);
    \begin{scope}[on background layer]
      \node[draw, dashed, black!60, fit=(ant)(tun)(adc), inner sep=7pt,
        label={[anchor=south west, black!70]north west:odbiornik RTL-SDR Blog V3}] {};
      \node[draw, dashed, black!60, fit=(mod)(pre)(ppm)(crc)(avr)(sbs), inner sep=7pt,
        label={[anchor=north west, black!70]south west:komputer: program dump1090}] {};
    \end{scope}
  \end{tikzpicture}
  \caption{Tor odbioru ramek ADS-B z~wykorzystaniem odbiornika RTL-SDR i~programu dump1090; opracowanie własne na podstawie~\cite{sun2021,dump1090}}
  \label{rys:tor_sdr}
\end{figure}

Częstotliwości \SI{1090}{\mega\hertz} odpowiada długość fali $c/f \approx \SI{27.5}{cm}$, gdzie $c$ jest prędkością światła, a~$f$ częstotliwością fali. Do odbioru sygnałów ADS-B można zatem użyć prostych anten: dipola półfalowego o~całkowitej długości około \SI{13.8}{cm} albo anteny ćwierćfalowej z~przeciwwagą, której element odbiorczy ma długość około \SI{6.9}{cm}~\cite{sun2021}.
Przy długim przewodzie antenowym sygnał o~tej częstotliwości jest silnie tłumiony. Stosuje się wtedy antenę aktywną ze wzmacniaczem niskoszumowym, zasilaną przez przewód antenowy za pośrednictwem układu \textit{bias tee}~\cite{sun2021}. Odbiornik RTL-SDR Blog~V3 ma wbudowany, włączany programowo układ \textit{bias tee} o~napięciu \SI{4.5}{V}~\cite{rtlsdrv3}.

Sygnały o~częstotliwości \SI{1090}{\mega\hertz} rozchodzą się w~zasięgu wzroku, dlatego maksymalny zasięg odbioru wynika z~krzywizny Ziemi~\cite{sun2021}:
\begin{equation}
  d = R_{\mathrm{Z}} \left( \arccos\frac{R_{\mathrm{Z}}}{R_{\mathrm{Z}} + h_{\mathrm{r}}}
    + \arccos\frac{R_{\mathrm{Z}}}{R_{\mathrm{Z}} + h_{\mathrm{t}}} \right),
  \label{eq:zasieg}
\end{equation}
\begin{gdzie}
  \item $d$ -- maksymalny zasięg odbioru, w~kilometrach,
  \item $R_{\mathrm{Z}}$ -- promień Ziemi, w~kilometrach,
  \item $h_{\mathrm{r}}$ -- wysokość anteny odbiornika nad poziomem morza, w~kilometrach,
  \item $h_{\mathrm{t}}$ -- wysokość statku powietrznego nad poziomem morza, w~kilometrach.
\end{gdzie}
Rzeczywisty zasięg jest mniejszy, ponieważ zależy także od mocy nadajnika i~kierunkowości anten~\cite{sun2021}. Wzór~\eqref{eq:zasieg} ma jednak istotną konsekwencję dla dekodowania lokalnego. Dla anteny umieszczonej na poziomie morza ($h_{\mathrm{r}} = 0$) i~$R_{\mathrm{Z}} = \SI{6371}{km}$ teoretyczny zasięg przekracza \SI{180}{NM} (\SI{333}{km}) dla statków powietrznych lecących wyżej niż około \SI{8.7}{km}, czyli na typowych wysokościach przelotowych. Dla takich statków powietrznych dekodowanie lokalne z~położeniem odbiornika jako punktem odniesienia może prowadzić do wyboru niewłaściwej strefy (podrozdział~\ref{sec:cpr}).


% ---------------------------------------------------------------------
\section{Pozycja i~czas z~odbiornika GNSS}
% =====================================================================


\label{sec:gnss}
Globalne systemy nawigacji satelitarnej (ang.~\textit{Global Navigation Satellite System}, GNSS) umożliwiają wyznaczenie położenia i~czasu w~dowolnym miejscu na Ziemi na podstawie sygnałów nadawanych przez satelity~\cite{hofmann2008}. Do systemów tych należą m.in. amerykański GPS, rosyjski GLONASS, europejski Galileo i~chiński BeiDou. Podstawowa konstelacja systemu GPS obejmuje 24~pozycje orbitalne rozmieszczone w~sześciu płaszczyznach orbit. Zapewnia to widoczność co najmniej czterech satelitów praktycznie z~każdego miejsca na Ziemi. W~praktyce utrzymuje się więcej satelitów, w~ostatnich latach około 31~\cite{gpsgov_space}.

Satelity GNSS mają na pokładzie zegary atomowe i~nadają sygnały zawierające m.in. czas nadania oraz parametry swoich orbit. Odbiornik mierzy czas propagacji sygnału każdego satelity i~wyznacza tzw. pseudoodległość, czyli odległość obarczoną błędem wynikającym z~niedokładności własnego zegara~\cite{hofmann2008}.
Dla satelity o~numerze $k$ pseudoodległość opisuje zależność~\cite{hofmann2008}
\begin{equation}
  p_k = \lVert \boldsymbol{s}_k - \boldsymbol{r} \rVert + c\,\Delta t_{\mathrm{odb}},
  \label{eq:pseudoodleglosc}
\end{equation}
\begin{gdzie}
  \item $p_k$ -- pseudoodległość do satelity $k$, w~metrach,
  \item $\boldsymbol{s}_k$ -- wektor położenia satelity $k$, w~metrach,
  \item $\boldsymbol{r}$ -- wektor położenia odbiornika, w~metrach,
  \item $c$ -- prędkość światła, w~metrach na sekundę,
  \item $\Delta t_{\mathrm{odb}}$ -- przesunięcie zegara odbiornika względem czasu systemu, w~sekundach.
\end{gdzie}
Niewiadomymi są trzy współrzędne odbiornika oraz przesunięcie jego zegara, dlatego do wyznaczenia pozycji potrzebne są pomiary z~co najmniej czterech satelitów. Ubocznym efektem obliczeń jest dokładna znajomość czasu, co pozwala wykorzystać odbiornik GNSS jako wzorzec czasu~\cite{hofmann2008}.

Odbiorniki GNSS wyrażają położenie w~układzie odniesienia WGS-84, w~którym kształt Ziemi przybliżono elipsoidą obrotową~\cite{wgs84}. Położenie opisuje się szerokością geograficzną $\varphi$, długością geograficzną $\lambda$ i~wysokością nad elipsoidą $h$. Często podaje się jednak wysokość nad poziomem morza $H$, liczoną względem geoidy, czyli powierzchni ekwipotencjalnej pola siły ciężkości zbliżonej do średniego poziomu mórz. Obie wysokości wiąże zależność~\cite{hofmann2008}
\begin{equation}
  h \approx H + N,
  \label{eq:wysokosci}
\end{equation}
\begin{gdzie}
  \item $N$ -- odstęp geoidy od elipsoidy, w~metrach.
\end{gdzie}
Odstęp ten osiąga wartości rzędu kilkudziesięciu metrów, dlatego w~obliczeniach geometrycznych wszystkie wysokości należy odnosić do tej samej powierzchni.

Odbiorniki GNSS przekazują wyniki najczęściej w~postaci tekstowych zdań standardu NMEA~0183, przesyłanych łączem szeregowym~\cite{nmea0183}. Każde zdanie rozpoczyna się znakiem \texttt{\$}, po którym następują dwuliterowy identyfikator nadawcy i~trzyliterowy typ zdania. Identyfikator \texttt{GP} oznacza dane z~systemu GPS, a~\texttt{GN} -- rozwiązanie wykorzystujące kilka systemów jednocześnie. Pola danych rozdziela się przecinkami. Zdanie kończą znak \texttt{*} i~dwucyfrowa szesnastkowa suma kontrolna, będąca wynikiem operacji XOR na kodach wszystkich znaków zawartych między \texttt{\$} a~\texttt{*}~\cite{gpsd_nmea}. Pola zdań GGA i~RMC zestawiono w~tabeli~\ref{tab:nmea}. Zdanie GGA zawiera m.in. wysokość nad poziomem morza i~odstęp geoidy od elipsoidy, co pozwala wyznaczyć wysokość nad elipsoidą ze wzoru~\eqref{eq:wysokosci}~\cite{gpsd_nmea}. Przykładowe zdania z~zaokrąglonymi współrzędnymi mają postać:
\begin{center}
  \small
  \begin{tabular}{@{}l@{}}
    \texttt{\$GPGGA,123519.00,5000.0000,N,02000.0000,E,1,08,0.9,250.0,M,40.0,M,,*64}\\
    \texttt{\$GPRMC,123519.00,A,5000.0000,N,02000.0000,E,0.0,0.0,061126,,*3B}
  \end{tabular}
\end{center}

\begin{table}[H]
  \centering
  \caption{Wybrane pola zdań GGA i~RMC standardu NMEA~0183 oraz ich wartości w~przykładowych zdaniach; opracowanie własne na podstawie~\cite{gpsd_nmea}}
  \label{tab:nmea}
  \begin{tabular}{@{}llll@{}}
    \toprule
    Zdanie & Pole & Zawartość & Przykład \\
    \midrule
    GGA & 1      & czas UTC (hhmmss.ss)                         & 12:35:19,00 \\
        & 2--3   & szerokość geograficzna (ddmm.mmmm, N/S)       & \ang{50.0} N \\
        & 4--5   & długość geograficzna (dddmm.mmmm, E/W)        & \ang{20.0} E \\
        & 6      & jakość wyznaczenia (0 -- brak, 1 -- GPS, 2 -- DGPS) & 1 \\
        & 7      & liczba wykorzystanych satelitów               & 8 \\
        & 8      & współczynnik HDOP                             & \num{0.9} \\
        & 9--10  & wysokość nad poziomem morza, w~metrach        & \num{250.0} \\
        & 11--12 & odstęp geoidy od elipsoidy, w~metrach         & \num{40.0} \\
    \midrule
    RMC & 1      & czas UTC (hhmmss.ss)                         & 12:35:19,00 \\
        & 2      & status (A -- dane poprawne, V -- niepoprawne)  & A \\
        & 3--6   & szerokość i~długość geograficzna              & jak w~GGA \\
        & 7      & prędkość względem ziemi, w~węzłach            & \num{0.0} \\
        & 8      & kąt drogi, w~stopniach                        & \num{0.0} \\
        & 9      & data (ddmmyy)                                 & 6.11.2026 \\
    \bottomrule
  \end{tabular}
\end{table}

Czas w~zdaniach NMEA jest podawany w~skali UTC, a~zdanie RMC zawiera także datę~\cite{gpsd_nmea}.
Czas systemu GPS różni się od czasu UTC o~całkowitą liczbę sekund przestępnych, którą odbiornik uwzględnia przy podawaniu czasu UTC~\cite{hofmann2008}.
Odbiornik GNSS może zatem służyć jako źródło czasu referencyjnego, niezależnego od zegara komputera. Sposób wykorzystania tego czasu do korekty znaczników czasu rejestrowanych danych opisano w~podrozdziale~\ref{sec:model}.


% ---------------------------------------------------------------------
\section{Środowisko ROS~2}
% =====================================================================
\label{sec:ros2}


Robot Operating System~2 (ROS~2) to zestaw bibliotek i~narzędzi do tworzenia oprogramowania robotów~\cite{ros2jazzy}. Mimo nazwy nie jest to system operacyjny, lecz platforma programistyczna, której podstawę stanowi warstwa pośrednicząca w~komunikacji między elementami aplikacji (ang.~\textit{middleware})~\cite{macenski2022}. ROS~2 działa w~systemie operacyjnym komputera, a~większość jego ekosystemu rozwija się wokół systemu Linux~\cite{ye2023realtime}. W~stosunku do pierwszej wersji ROS~2 zaprojektowano od nowa, m.in. z~myślą o~systemach wbudowanych, pracy w~czasie rzeczywistym i~komunikacji w~sieciach o~ograniczonej niezawodności~\cite{macenski2022}.

Oprogramowanie ROS~2 jest wydawane w~dystrybucjach. W~pracy wykorzystano dystrybucję Jazzy Jalisco, wydaną w~maju 2024~roku jako wersja o~wydłużonym okresie wsparcia, trwającym do maja 2029~roku. Najwyższy poziom wsparcia ma ona w~systemach Ubuntu~24.04, zarówno dla procesorów amd64, jak i~arm64, oraz Windows~10~\cite{ros2jazzyrelease}.

Aplikacja w~ROS~2 składa się z~węzłów (ang.~\textit{node}), czyli niezależnych jednostek programu realizujących wydzielone zadania~\cite{ros2jazzy}. Węzły wymieniają dane głównie przez tematy (ang.~\textit{topic}) w~modelu publikacji i~subskrypcji. Węzeł publikujący wysyła komunikaty (ang.~\textit{message}) do nazwanego tematu, nie wiedząc, kto je odbiera. Węzły subskrybujące dany temat otrzymują publikowane w~nim komunikaty, przy czym gwarancja ich dostarczenia zależy od ustawień jakości usługi, omówionych w~dalszej części podrozdziału~\cite{ros2jazzy}. Jeden temat może mieć wielu nadawców i~wielu odbiorców, co przedstawiono na rysunku~\ref{rys:ros2_graf}.

\begin{figure}[H]
  \centering
  \begin{tikzpicture}[font=\footnotesize, >=Stealth,
      wezel/.style={draw, thick, ellipse, fill=akcent!8, align=center, minimum height=0.95cm, inner sep=2pt},
      temat/.style={draw, thick, rectangle, fill=black!5, align=center, minimum height=0.8cm, inner sep=5pt},
      plik/.style={draw, thick, cylinder, shape border rotate=90, aspect=0.25, fill=black!5, align=center,
        minimum height=1.1cm, minimum width=1.3cm}]
    \node[wezel] (a) {węzeł A\\ (nadawca)};
    \node[temat, right=1.6cm of a] (t) {\texttt{/temat}};
    \node[wezel, right=2.0cm of t, yshift=1.25cm] (b) {węzeł B};
    \node[wezel, right=2.0cm of t] (c) {węzeł C};
    \node[wezel, right=2.0cm of t, yshift=-1.25cm] (r) {rejestrator\\ rosbag2};
    \node[plik, right=1.1cm of r] (f) {plik\\ MCAP};
    \draw[->, thick] (a) -- node[above] {publikacja} (t);
    \draw[->, thick] (t.east) -- node[sloped, above, pos=0.5] {subskrypcja} (b.west);
    \draw[->, thick] (t.east) -- (c.west);
    \draw[->, thick] (t.east) -- (r.west);
    \draw[->, thick] (r) -- node[above] {zapis} (f);
  \end{tikzpicture}
  \caption{Komunikacja węzłów ROS~2 przez temat w~modelu publikacja--subskrypcja; węzły oznaczono elipsami, a~temat -- prostokątem; rejestrator rosbag2 i~format MCAP omówiono w~dalszej części podrozdziału; opracowanie własne}
  \label{rys:ros2_graf}
\end{figure}

Komunikację między węzłami realizuje warstwa pośrednia zgodna ze standardem DDS (ang.~\textit{Data Distribution Service}). Odpowiada ona za automatyczne wykrywanie węzłów w~sieci, bez centralnego serwera, oraz za przesyłanie danych~\cite{macenski2022}. Domyślną implementacją DDS w~dystrybucji Jazzy jest eProsima Fast DDS~\cite{ros2jazzy}. Dzięki rozproszonemu wykrywaniu węzły uruchomione na różnych komputerach w~tej samej sieci mogą wymieniać komunikaty bez dodatkowej konfiguracji, o~ile mają ten sam identyfikator domeny, ustawiany zmienną środowiskową \texttt{ROS\_DOMAIN\_ID}. Przypisanie różnych identyfikatorów pozwala oddzielić od siebie grupy węzłów pracujące w~jednej sieci~\cite{ros2jazzy}.

Sposób przesyłania komunikatów określają ustawienia jakości usługi (ang.~\textit{quality of service}, QoS)~\cite{ros2jazzy}. Należą do nich m.in. niezawodność dostarczania oraz historia. W~trybie niezawodnym utracone komunikaty są wysyłane ponownie, a~w~trybie \textit{best effort} nie są. Historia określa liczbę komunikatów przechowywanych w~kolejce, nazywaną głębokością kolejki. Gdy kolejka jest pełna, najstarsze komunikaty są usuwane. Ustawienia nadawcy i~odbiorcy muszą być zgodne, w~przeciwnym razie połączenie między nimi nie powstaje~\cite{ros2jazzy}.

Opóźnienie przesyłania komunikatu między węzłami nie jest stałe. Zależy ono m.in. od implementacji DDS, rozmiaru komunikatu i~ustawień QoS. Dla małych komunikatów przesyłanych między węzłami działającymi na jednym komputerze jest ono rzędu ułamka milisekundy~\cite{ye2023realtime,kronauer2021latency}. Narzut samego ROS~2 w~stosunku do bezpośredniego użycia DDS może sięgać 50\,\% dla małych komunikatów, a~tryb niezawodny zwiększa opóźnienie o~kilkanaście procent w~porównaniu z~trybem \textit{best effort}~\cite{kronauer2021latency}.

Strukturę komunikatu opisuje plik interfejsu z~rozszerzeniem \texttt{.msg}, w~którym wymienia się pola i~ich typy~\cite{ros2jazzy}. Na tej podstawie narzędzia ROS~2 generują kod obsługujący komunikat w~językach C++ i~Python. Dostępne są standardowe typy komunikatów, np. \texttt{sensor\_msgs/NavSatFix}, przeznaczony do przekazywania pozycji z~odbiornika GNSS, oraz \texttt{sensor\_msgs/TimeReference}, przeznaczony do przekazywania czasu z~zewnętrznego źródła~\cite{ros2jazzy}. Gdy żaden typ standardowy nie odpowiada potrzebom, można zdefiniować własny typ komunikatu. Wiele typów komunikatów zawiera nagłówek ze znacznikiem czasu i~nazwą układu współrzędnych, do którego odnoszą się dane~\cite{ros2jazzy}. Znacznik czasu nadaje węzeł publikujący. Domyślnie jest to czas zegara systemowego komputera, na którym działa węzeł, a~po ustawieniu parametru \texttt{use\_sim\_time} -- czas publikowany w~temacie \texttt{/clock}~\cite{ros2jazzy}. Dokładność znaczników czasu zależy zatem od synchronizacji zegara komputera, co uwzględniono w~modelu opisanym w~podrozdziale~\ref{sec:model}.

Do rejestracji danych służy narzędzie rosbag2, które zapisuje komunikaty wybranych tematów do pliku wraz ze znacznikami czasu i~pozwala je później odtworzyć~\cite{ros2jazzy}. Odtwarzane komunikaty są publikowane w~tych samych tematach, dlatego węzły przetwarzające dane działają tak samo jak podczas pracy na żywo. Podczas odtwarzania można publikować również zegar odpowiadający czasowi nagrania. Węzły z~ustawionym parametrem \texttt{use\_sim\_time} posługują się wtedy czasem zapisu zamiast czasem bieżącym~\cite{ros2jazzy}. Od dystrybucji Iron, wydanej w~maju 2023~roku, domyślnym formatem zapisu jest MCAP, który zastąpił stosowany wcześniej format SQLite3~\cite{ros2iron}. Plik MCAP zawiera definicje zapisanych komunikatów, dzięki czemu można go odczytać także bez zainstalowanych definicji typów komunikatów~\cite{ros2iron,mcap}.


% ---------------------------------------------------------------------
\section{Istniejące rozwiązania}
% =====================================================================
\label{sec:istniejace}


Odbiór sygnałów ADS-B za pomocą tanich odbiorników SDR jest powszechny, a~budowa własnego odbiornika jest popularnym projektem hobbystycznym. Typowy odbiornik składa się z~anteny, odbiornika RTL-SDR i~dekodera, np. programu dump1090~\cite{sun2021}, uruchamianego często na komputerze jednopłytkowym Raspberry Pi~\cite{fr24build,piaware}. Odebrane statki powietrzne wyświetla się zwykle na mapie w~przeglądarce internetowej, a~zdekodowane dane są dostępne przez porty TCP programu dump1090~\cite{dump1090}. Takie odbiorniki służą przede wszystkim do obserwacji ruchu lotniczego w~okolicy. Dane nie są w~nich udostępniane jako komunikaty środowiska ROS~2, a~w~podstawowej konfiguracji nie są też zapisywane do plików.

Odbiorniki budowane przez hobbystów często przekazują dane do serwisów śledzenia lotów. Serwis Flightradar24 udostępnia instrukcję budowy odbiornika złożonego z~komputera Raspberry Pi, odbiornika ADS-B w~postaci przystawki USB oraz anteny. Na komputerze instaluje się przygotowany przez serwis obraz systemu Pi24, oparty na systemie Raspbian Lite, a~odbiornik automatycznie przesyła odebrane dane do serwisu. Według serwisu zasięg takiego odbiornika wynosi od 200 do 400~mil, czyli od około 320 do \SI{640}{km}. W~zamian za udostępnianie danych użytkownik otrzymuje bezpłatny dostęp do płatnego planu subskrypcji~\cite{fr24build}. Podobnie działa oprogramowanie PiAware, które na komputerze Raspberry Pi dekoduje wiadomości ADS-B, wyświetla je na lokalnej mapie i~przesyła do serwisu FlightAware. Jego użytkownicy otrzymują bezpłatne konto z~rozszerzonym dostępem do danych serwisu~\cite{piaware}.

Przykładem rozwiązania naukowego jest OpenSky Network, czyli sieć odbiorników naziemnych, która gromadzi surowe wiadomości Mode~S i~ADS-B w~centralnej bazie danych i~udostępnia je do celów badawczych~\cite{schafer2014}. Wszystkie opisane sieci są nastawione na gromadzenie danych z~wielu odbiorników w~jednym miejscu. Do przesyłania danych wymagają połączenia z~internetem, a~dane udostępniają przez własne serwisy.

Dane ADS-B wykorzystuje się także w~badaniach nad wizyjnym wykrywaniem statków powietrznych. Zhang i~in.~\cite{zhang2022sta} zaproponowali sieć neuronową do segmentacji statków powietrznych w~nagraniach wideo z~lotniska. Położenie odczytane z~wiadomości ADS-B służy w~niej jako informacja wstępna, która ułatwia śledzenie obiektów w~długich sekwencjach. W~zbiorze danych TartanAviation~\cite{patrikar2025tartan} zarejestrowano obrazy z~kamer zamontowanych na stałe w~obrębie lotnisk wraz z~danymi ADS-B z~własnych odbiorników. Położenia statków powietrznych interpolowano do częstotliwości rejestracji obrazu, wynoszącej \SI{24}{\hertz}, i~rzutowano na obraz za pomocą skalibrowanego modelu kamer. Rzutowane położenia posłużyły do odrzucania błędnych wykryć i~wskazania przybliżonego położenia statku powietrznego na obrazie, a~etykiety sprawdzono następnie ręcznie~\cite{patrikar2025tartan}. Przykłady te pokazują, że do automatycznego etykietowania obrazów potrzebne są zsynchronizowany czas, znane położenie odbiornika i~kamery oraz wiarygodne położenia statków powietrznych.

\begin{table}[htbp]
  \centering
  \caption{Porównanie odbiorników ADS-B o~różnym przeznaczeniu; opracowanie własne na podstawie~\cite{fr24build,piaware,schafer2014,dump1090}}
  \label{tab:rozwiazania}
  \small
  \begin{tabular}{@{}>{\raggedright\arraybackslash}p{2.8cm}>{\raggedright\arraybackslash}p{3.9cm}>{\raggedright\arraybackslash}p{3.9cm}>{\raggedright\arraybackslash}p{3.9cm}@{}}
    \toprule
    Cecha & Sieci odbiorników (Flightradar24, FlightAware, OpenSky) & Odbiornik hobbystyczny z~programem dump1090 & Stacja bazowa (niniejsza praca) \\
    \midrule
    Przeznaczenie & gromadzenie danych z~wielu odbiorników i~ich udostępnianie & obserwacja ruchu lotniczego w~okolicy odbiornika & źródło danych referencyjnych w~terenie \\
    \addlinespace
    Połączenie z~internetem & wymagane do przesyłania danych & niewymagane & niewymagane \\
    \addlinespace
    Postać udostępnianych danych & serwis internetowy lub baza danych & mapa w~przeglądarce, porty TCP & komunikaty ROS~2 w~czasie rzeczywistym \\
    \addlinespace
    Zapis danych & centralnie, w~serwisie & brak w~podstawowej konfiguracji & lokalnie, pliki rosbag2 w~formacie MCAP \\
    \bottomrule
  \end{tabular}
\end{table}

Cechy omówionych odbiorników zestawiono w~tabeli~\ref{tab:rozwiazania}. W~przeanalizowanych rozwiązaniach nie znaleziono przenośnego systemu, który w~warunkach polowych, bez dostępu do internetu, wyznaczałby położenia statków powietrznych względem własnej pozycji ustalonej przez odbiornik GNSS i~udostępniał je w~czasie rzeczywistym jako komunikaty środowiska ROS~2. Wymagania stawiane takiemu systemowi sformułowano w~podrozdziale~\ref{sec:wymagania_referencja}.
                 % 2. ADS-B, CPR, SDR, GNSS, ROS 2, istniejące rozwiązania
% =====================================================================
%  3. PROJEKT STACJI BAZOWEJ  (wersja 2026-10-29)
%  Objętość: ok. 9 stron. Model (35_Model.tex) dołączony na końcu.
% =====================================================================
\chapter{Projekt stacji bazowej}
\label{chap:projekt}

% 2-3 zdania: co jest w tym rozdziale.

% ---------------------------------------------------------------------
\section{Wymagania}
\label{sec:wymagania}

\subsection{Wymagania funkcjonalne}
\label{sec:wymagania_funkcjonalne}
% Odbiór ramek ADS-B, odczyt pozycji i czasu GPS, dekodowanie pozycji
% w czasie rzeczywistym, publikacja jako komunikaty ROS 2, rejestracja,
% podgląd stanu na wyświetlaczu.

\subsection{Wymagania pozafunkcjonalne}
\label{sec:wymagania_pozafunkcjonalne}
% Wielogodzinna praca z powerbanka, przenośność, start bez obsługi
% (autostart), odporność na nagłą utratę zasilania, praca bez internetu.

\subsection{Wymagania dotyczące danych referencyjnych}
\label{sec:wymagania_referencja}
% Znaczniki czasu UTC o znanej dokładności, znana pozycja stacji,
% bezstratny zapis surowych danych, jednoznaczne pozycje samolotów,
% kierunek do samolotu (azymut, elewacja).

% ---------------------------------------------------------------------
\section{Architektura systemu}
\label{sec:architektura}
% (ok. 2 str.) Podział na warstwę sprzętową i programową; podział
% demodulacja (dump1090) / dekodowanie (decoder_node); diagram przepływu
% danych; scenariusz pracy w terenie (włączenie -> nagrywanie ->
% wyłączenie przyciskiem -> skopiowanie danych).
% ======================================================================
% rys_sprzet.tex – schemat blokowy warstwy sprzętowej stacji bazowej
% Wstawić w rozdziale 3 (Architektura) poleceniem: \input{tex/rys_sprzet}
% Wymaga w preambule: tikz + biblioteki positioning, arrows.meta, calc,
%                     decorations.pathmorphing; siunitx; kolor akcent.
% Kod TikZ przygotowano z wykorzystaniem Claude (Anthropic) – zob. deklarację AI.
% ======================================================================

Warstwę sprzętową stacji bazowej przedstawiono na rysunku~\ref{rys:sprzet}. Centralnym elementem jest komputer Raspberry Pi~4B, do którego przez port USB dołączono odbiornik SDR, a przez interfejs UART moduł odbiornika GPS. Wyświetlacz TFT, dioda stanu i~przycisk zasilania są połączone bezpośrednio z~wyprowadzeniami GPIO.

\begin{figure}[!htb]
\centering
\begin{tikzpicture}[
    font=\footnotesize,
    >={Stealth[length=2mm]},
    blok/.style={draw, thick, fill=akcent!10, align=center,
                 minimum height=9mm, inner sep=3pt, text width=#1},
    blok/.default=27mm,
    zewn/.style={align=center, text width=34mm, font=\footnotesize},
    dane/.style={->, thick},
    zasil/.style={->, line width=2pt, draw=black!50},
    opcja/.style={<->, thick, dashed},
    radio/.style={->, thick, decorate,
                  decoration={snake, amplitude=0.6mm, segment length=2.5mm,
                              post length=2mm}},
    etyk/.style={font=\scriptsize, align=center, fill=white, inner sep=1pt},
]
% ---------------- Raspberry Pi (środek) ----------------
\node[draw, very thick, fill=akcent!20, align=center,
      minimum width=44mm, minimum height=44mm] (rpi) at (0,0)
     {\textbf{Raspberry Pi 4B}\\[2mm]
      \scriptsize Ubuntu, Docker,\\ \scriptsize ROS~2 Jazzy,\\
      \scriptsize dump1090};

% ---------------- lewa strona: odbiór radiowy ----------------
\node[blok] (sdr) at (-6.3, 1.3) {Odbiornik SDR\\RTL-SDR Blog V3};
\node[blok] (gps) at (-6.3,-1.3) {Moduł odbiornika GPS};
\node[zewn] (adsb) at (-6.3, 3.6)
     {statki powietrzne\\(ADS-B, \SI{1090}{\mega\hertz})};
\node[zewn] (sat) at (-6.3,-3.6)
     {satelity GPS\\(L1, \SI{1575.42}{\mega\hertz})};
\draw[radio] (adsb) -- node[etyk, right=1mm, fill=none]
     {antena dipolowa $\lambda/4$} (sdr);
\draw[radio] (sat) -- node[etyk, right=1mm, fill=none]
     {antena GPS} (gps);
\draw[dane] (sdr.east) -- node[etyk, above=1pt, fill=none] {USB~2.0}
     (sdr.east -| rpi.west);
\draw[dane] (gps.east) -- node[etyk, above=1pt, fill=none]
     {UART} node[etyk, below=1pt, fill=none] {\SI{9600}{bit/s}}
     (gps.east -| rpi.west);

% ---------------- prawa strona: interfejs użytkownika ----------------
\node[blok] (tft) at (6.3, 1.5) {Wyświetlacz TFT\\ILI9341, 240$\times$320};
\node[blok] (led) at (6.3, 0.0) {Dioda LED stanu};
\node[blok] (btn) at (6.3,-1.5) {Przycisk zasilania};
\draw[dane] (tft.west -| rpi.east) -- node[etyk, above=1pt, fill=none] {SPI0}
     node[etyk, below=1pt, fill=none] {GPIO24/25} (tft.west);
\draw[dane] (led.west -| rpi.east) -- node[etyk, above=1pt, fill=none]
     {GPIO17} (led.west);
\draw[dane] (btn.west) -- node[etyk, above=1pt, fill=none] {GPIO3}
     (btn.west -| rpi.east);

% ---------------- dół: zasilanie i pamięć ----------------
\node[blok=31mm] (pb)  at (-1.9,-4.0) {Powerbank\\LYNXER RPP-272};
\node[blok=31mm] (usb) at ( 1.9,-4.0) {Pamięć USB \SI{64}{GB}\\(system, pliki MCAP)};
\draw[zasil] (pb.north) -- node[etyk, left=1pt, fill=none]
     {USB-C, \SI{5}{V}} (pb.north |- rpi.south);
\draw[<->, thick] (usb.north) -- node[etyk, right=1pt, fill=none] {USB}
     (usb.north |- rpi.south);

% ---------------- góra: opcjonalny komputer ----------------
\node[blok=36mm, fill=white, dashed] (pc) at (0, 3.9)
     {Komputer do analizy\\(opcjonalnie)};
\draw[opcja] (pc.south) -- node[etyk, right=1pt, fill=none]
     {Wi-Fi, DDS / SSH} (rpi.north);
\end{tikzpicture}
\caption{Schemat blokowy warstwy sprzętowej stacji bazowej. Strzałki ciągłe
oznaczają kierunek przepływu danych, gruba szara strzałka -- zasilanie,
strzałki faliste -- sygnały radiowe, linia przerywana -- połączenie
opcjonalne. Przy strzałkach podano interfejsy i numery wyprowadzeń GPIO
(numeracja BCM).}
\label{rys:sprzet}
\end{figure}
\newpage
% ---------------------------------------------------------------------
\section{Tematy i~komunikaty}
\label{sec:tematy}
% (ok. 1 str.) Lista tematów (tabela) + definicja komunikatu AircraftState
% (pola, jednostki SI, NaN = brak wartości, stałe position_quality),
% uzasadnienie własnego typu zamiast JSON w std_msgs/String.

% ======================================================================
% rys_graf_ros.tex – graf obliczeniowy ROS 2 stacji bazowej
% (węzły, tematy, typy komunikatów, źródła danych w systemie gospodarza)
% Wstawić w rozdziale 3 (Architektura) poleceniem: \input{tex/rys_graf_ros}
% Wymaga w preambule: tikz + biblioteki positioning, arrows.meta, fit,
%                     backgrounds, calc; kolor akcent.
% Nazwy węzłów i tematów zgodne z repozytorium (commit 33bb6ff).
% Kod TikZ przygotowano z wykorzystaniem Claude (Anthropic) – zob. deklarację AI.
% ======================================================================
\begin{figure}[!htb]
\centering
\begin{tikzpicture}[
    font=\footnotesize,
    >={Stealth[length=1.8mm]},
    wezel/.style={draw, thick, rounded corners=3mm, fill=akcent!20,
                  align=center, minimum height=8mm, text width=27mm,
                  font=\scriptsize\ttfamily},
    temat/.style={draw, thick, fill=white, align=center,
                  minimum height=8mm, text width=32mm, inner sep=2pt},
    host/.style={draw, thick, fill=black!7, align=center,
                 minimum height=8mm, text width=23mm, font=\scriptsize},
    pol/.style={->, thick},
    etyk/.style={font=\scriptsize, align=center, inner sep=1pt},
]
% makro: nazwa tematu (czcionka maszynowa) + typ komunikatu (kursywa)
\newcommand{\tm}[2]{{\scriptsize\ttfamily #1}\\[-0.6mm]{\scriptsize\itshape #2}}

% ---------------- system gospodarza (źródła danych) ----------------
\node[host] (d1090) at (0, 3.65)
     {\textbf{dump1090}\\(demodulacja)\\TCP 30003: SBS\\TCP 30002: AVR};
\node[host] (uart)  at (0, 1.25) {\texttt{/dev/ttyAMA0}\\(UART, GPS)};
\node[host] (hwmon) at (0,-1.15) {\texttt{/sys/class/hwmon}\\(czujniki RPi)};
\draw[pol] ($(d1090.north)+(0,0.6)$) node[above, etyk] {z RTL-SDR (USB)}
     -- (d1090.north);

% ---------------- węzły publikujące ----------------
\node[wezel] (sdr)  at (3.5, 3.65) {sdr\_radar\_node};
\node[wezel] (gpsn) at (3.5, 1.25) {gps\_node};
\node[wezel] (pwr)  at (3.5,-1.15) {power\_monitor\_node};
\draw[pol] ([yshift=2mm]d1090.east)  -- ([yshift=2mm]sdr.west);
\draw[pol] ([yshift=-2mm]d1090.east) -- ([yshift=-2mm]sdr.west);
\draw[pol] (uart)  -- (gpsn);
\draw[pol] (hwmon) -- (pwr);

% ---------------- tematy ----------------
\node[temat] (raw)    at (6.95, 4.25) {\tm{/aircraft/raw}{std\_msgs/String}};
\node[temat] (frames) at (6.95, 3.05) {\tm{/aircraft/frames}{std\_msgs/String}};
\node[temat] (fix)    at (6.95, 1.85) {\tm{/gps/fix}{sensor\_msgs/NavSatFix}};
\node[temat] (tref)   at (6.95, 0.65) {\tm{/gps/time\_ref}{sensor\_msgs/TimeReference}};
\node[temat] (uv)     at (6.95,-0.55) {\tm{/system/undervoltage}{std\_msgs/Bool}};
\node[temat] (temp)   at (6.95,-1.75) {\tm{/system/cpu\_temp}{sensor\_msgs/Temperature}};

\draw[pol] (sdr.east)  -- (raw.west);
\draw[pol] (sdr.east)  -- (frames.west);
\draw[pol] (gpsn.east) -- (fix.west);
\draw[pol] (gpsn.east) -- (tref.west);
\draw[pol] (pwr.east)  -- (uv.west);
\draw[pol] (pwr.east)  -- (temp.west);

% ---------------- przetwarzanie i prezentacja ----------------
\node[wezel] (dec) at (10.45, 2.45) {decoder\_node};
\node[temat] (state) at (10.45, 0.75)
     {\tm{/aircraft/state}{aircraft\_interfaces/\\[-0.6mm]AircraftState}};
\node[wezel, text width=19mm] (disp) at (13.65,-0.55) {display\_node};
\node[etyk, below=1mm of disp.south east, anchor=north east] {$\rightarrow$ TFT (SPI)};

\draw[pol] (frames.east) -- ([yshift=1.5mm]dec.west);
\draw[pol] (fix.east)    -- ([yshift=-1.5mm]dec.west);
\draw[pol] (dec.south)   -- (state.north);
\draw[pol] (fix.east)   -| ([xshift=3mm]disp.north);
\draw[pol] (state.east) -| ([xshift=-3mm]disp.north);
\draw[pol] (uv.east)    -- (disp.west);
\draw[pol] (temp.east)  -| ([xshift=-6mm]disp.south);

% ---------------- moduł adsb_decoder (nie jest węzłem) ----------------
\node[draw, thick, dashed, fill=white, align=center, text width=29mm,
      minimum height=8mm, inner sep=2pt] (mod) at (10.95, 4.25)
     {{\scriptsize\ttfamily adsb\_decoder.py}\\[-0.6mm]
      {\scriptsize\itshape moduł Pythona (nie węzeł)}};
\draw[->, dashed, thick, >={Straight Barb[length=1.8mm]}]
     (mod.south |- dec.north) -- node[etyk, right=1mm] {«import»} (mod.south);

% ---------------- rejestracja rosbag2 ----------------
% ramka w kształcie litery L: kolumna tematów + temat /aircraft/state
% (bez węzła decoder_node, który nie jest tematem)
\coordinate (gL) at ($(raw.west)+(-2.5mm,0)$);
\coordinate (gR) at ($(raw.east)+(2.5mm,0)$);
\coordinate (gT) at ($(raw.north)+(0,2.5mm)$);
\coordinate (gB) at ($(temp.south)+(0,-2.5mm)$);
\coordinate (sT) at ($(state.north)+(0,2mm)$);
\coordinate (sB) at ($(state.south)+(0,-2mm)$);
\coordinate (sR) at ($(state.east)+(2mm,0)$);
\draw[dashed, rounded corners=2mm]
     (gL |- gT) -- (gR |- gT) -- (gR |- sT) -- (sR |- sT) -- (sR |- sB)
     -- (gR |- sB) -- (gR |- gB) -- (gL |- gB) -- cycle;
\node[wezel, text width=30mm] (bag) at (6.95,-3.4) {rosbag2\_recorder};
\node[etyk, right=2mm of bag, align=left] (mcap)
     {$\rightarrow$ pliki \texttt{.mcap}\\\phantom{$\rightarrow$} (co \SI{10}{min})};
\draw[->, line width=1.4pt] (gB -| bag.north) -- node[etyk, left=1mm]
     {wszystkie tematy z~ramki} (bag.north);

% ---------------- granice: Ubuntu (poza kontenerem) i kontener ----------------
\begin{scope}[on background layer]
  \node[draw, thick, densely dotted, rounded corners=2mm,
        fit=(sdr)(pwr)(raw)(temp)(state)(disp)(mod)(bag)(mcap), inner sep=3mm,
        label={[etyk, anchor=south east]north east:kontener Docker (ROS~2 Jazzy)}]
        (kont) {};
  \node[draw, thick, densely dotted, rounded corners=2mm,
        fit=(d1090)(hwmon), inner sep=1.5mm,
        label={[etyk, anchor=north, text width=25mm]south:Ubuntu na RPi\\(poza kontenerem)}] (gosp) {};
\end{scope}
\end{tikzpicture}
\caption{Graf obliczeniowy ROS~2 stacji bazowej. Prostokąty zaokrąglone
oznaczają węzły, prostokąty o~ostrych narożnikach -- tematy wraz z~typem
komunikatu (kursywa), szare prostokąty -- źródła danych w~systemie Ubuntu
uruchomionym bezpośrednio na Raspberry Pi, poza kontenerem. Strzałka ciągła
prowadzi od wydawcy do subskrybenta. Prostokąt przerywany oznacza moduł Pythona
\texttt{adsb\_decoder.py}, który nie jest węzłem ROS~2, a~przerywana strzałka
«import» -- jego wykorzystanie przez węzeł \texttt{decoder\_node}. Ramka
przerywana obejmuje tematy rejestrowane przez rosbag2, ramki kropkowane --
elementy uruchamiane w~kontenerze Docker i~poza nim.}
\label{rys:graf_ros}
\end{figure}

Przepływ danych w~środowisku ROS~2 pokazano na rysunku~\ref{rys:graf_ros}. Zgodnie z~typową praktyką w~ROS~2 każdy węzeł ma jedno, wyraźnie określone zadanie, a~dane między węzłami przekazywane są wyłącznie przez tematy. Węzeł \texttt{sdr\_radar\_node} odbiera od programu dump1090 zdekodowane wiadomości SBS oraz surowe ramki w~formacie AVR i~publikuje je w~dwóch osobnych tematach. Węzeł \texttt{gps\_node} odczytuje zdania NMEA z~odbiornika GPS, a~następnie publikuje pozycję stacji bazowej oraz czas GPS, każde w~osobnym temacie. Węzeł \texttt{decoder\_node} subskrybuje temat z~ramkami i~temat z~pozycją stacji, wyznacza na ich podstawie stan statków powietrznych względem stacji bazowej i~publikuje go w~temacie \texttt{/aircraft/state}. Do dekodowania ramek wykorzystuje moduł \texttt{adsb\_decoder.py}, który nie jest węzłem. Węzeł \texttt{power\_monitor\_node} odczytuje z~systemu temperaturę procesora oraz informację o~zbyt niskim napięciu zasilania i~publikuje je w~tematach \texttt{/system/cpu\_temp} i~\texttt{/system/undervoltage}. Węzeł \texttt{display\_node} subskrybuje te dwa tematy, a~także temat ze stanem statków powietrznych i~temat z~pozycją stacji, po czym wyświetla dane na ekranie TFT za pomocą biblioteki luma.lcd. Wszystkie tematy są zapisywane do plików MCAP przez rejestrator rosbag2.

\begin{table}[!htb]
  \centering
  \caption{Tematy ROS~2 publikowane przez stację bazową}
  \label{tab:tematy}
  \footnotesize
  \begin{tabular}{>{\raggedright}p{3.3cm} >{\raggedright}p{4.0cm} p{1.7cm} >{\raggedright\arraybackslash}p{4.9cm}}
    \toprule
    Temat & Typ komunikatu & Częstotl.\ [Hz] & Zawartość \\
    \midrule
    \texttt{/aircraft/raw}        & \texttt{std\_msgs/String}         & 0--50   & linia SBS z~dekodera dump1090 \\
    \texttt{/aircraft/frames}     & \texttt{std\_msgs/String}         & 30--70  & surowa ramka Mode~S w~formacie AVR \\
    \texttt{/aircraft/state}      & \texttt{aircraft\_interfaces/ AircraftState} & zmienna & stan statku powietrznego i~geometria względem stacji \\
    \texttt{/gps/fix}             & \texttt{sensor\_msgs/NavSatFix}   & 1       & pozycja stacji \\
    \texttt{/gps/time\_ref}       & \texttt{sensor\_msgs/ TimeReference} & 1    & czas UTC z~odbiornika GPS \\
    \texttt{/system/ undervoltage} & \texttt{std\_msgs/Bool}           & 1       & sygnalizacja zbyt niskiego napięcia zasilania \\
    \texttt{/system/cpu\_temp}    & \texttt{sensor\_msgs/ Temperature} & 1       & temperatura procesora \\
    \bottomrule
  \end{tabular}
\end{table}

% ---------------------------------------------------------------------
%  3.4. Model wyznaczania danych referencyjnych (osobny plik, zaczyna
%  się od \section). Umieszczony PO architekturze i tematach, bo model
%  korzysta z pojęć tam wprowadzonych (/gps/time_ref, znaczniki czasu).
% ---------------------------------------------------------------------
% ============================================================
%  35_Model.tex  --  Model opisowy i matematyczny wyznaczania
%                    danych referencyjnych (wersja uproszczona)
%  Wstawiany do rozdzialu 3 poleceniem:  % ============================================================
%  35_Model.tex  --  Model opisowy i matematyczny wyznaczania
%                    danych referencyjnych (wersja uproszczona)
%  Wstawiany do rozdzialu 3 poleceniem:  \input{tex/35_Model}
%  Wymaga w preambule: amsmath, siunitx, enumitem
% ============================================================

% Objaśnienia symboli pod wzorem ("gdzie:"): symbol w kolumnie od
% wcięcia akapitu, myślniki zawsze w jednej pionowej linii
\newenvironment{objasnienia}{%
  \noindent gdzie:
  \begin{description}[font=\normalfont, labelindent=\parindent,
    labelwidth=1.2cm, labelsep=0.2cm, leftmargin=!, align=left,
    itemsep=0pt, topsep=2pt]}{\end{description}}

\section{Model wyznaczania danych referencyjnych}
\label{sec:model}

Zarejestrowane przez stację dane mają posłużyć jako dane referencyjne, czyli
mają pozwolić stwierdzić, gdzie znajdował się dany statek powietrzny
w określonej chwili czasu. % referencja do podrozdziału.
W tym celu potrzebny jest model, który łączy trzy źródła informacji:
komunikaty ADS-B odbierane drogą radiową, pozycję stacji z odbiornika GPS oraz
zegar systemowy komputera, którym oznaczane są wszystkie zapisane wiadomości.
W niniejszym podrozdziale model ten przedstawiono najpierw w postaci słownej
(model opisowy), a następnie w postaci wzorów (model matematyczny). Na końcu
omówiono przyjęte uproszczenia i podano przykład obliczeniowy.

\subsection{Model opisowy}
\label{sec:model_opisowy}

Statek powietrzny wyposażony w transponder ADS-B kilka razy na sekundę
rozgłasza na częstotliwości \SI{1090}{MHz} komunikaty zawierające m.in.\ jego
unikalny adres ICAO, pozycję (szerokość i długość geograficzną) wyznaczoną
przez pokładowy odbiornik GNSS, wysokość barometryczną oraz prędkość
względem ziemi~\cite{sun2021}. Stacja odbiera część tych komunikatów --
niektóre giną z powodu zakłóceń, zbyt słabego sygnału lub przesłonięcia
anteny. Odebrane komunikaty są zapisywane w postaci tekstowej (format SBS)
razem ze znacznikiem czasu nadawanym przez komputer w chwili odbioru.

Komunikat ADS-B mówi, gdzie samolot jest, ale nie mówi, gdzie jest
względem stacji. Do tego potrzebna jest pozycja samej stacji, którą
podaje odbiornik GPS w komunikacie NMEA \texttt{GGA}~\cite{nmea0183}.
Znając obie pozycje, można obliczyć odległość od stacji do samolotu.

Drugim problemem jest czas. Komputer Raspberry Pi nie ma podtrzymywanego
bateryjnie zegara, a w terenie nie ma dostępu do internetu, więc jego zegar
może się mylić o nieznaną wartość. Ponieważ wszystkie zapisane wiadomości są
oznaczone właśnie tym zegarem, błąd zegara przenosi się na wszystkie dane.
Odbiornik GPS co sekundę przesyła jednak dokładny czas UTC (komunikat
\texttt{RMC}). Komputer zapisuje go razem z własnym wskazaniem zegara, dzięki
czemu po zakończeniu pomiarów można zmierzyć, o ile zegar się mylił,
i poprawić czas wszystkich wiadomości.

Granice modelowanego układu wyznacza sama stacja (antena, odbiornik RTL-SDR,
odbiornik GPS i komputer). Samoloty i satelity GPS należą do otoczenia,
a punktami styku są: łącze radiowe \SI{1090}{MHz}, łącze szeregowe
odbiornika GPS oraz plik zapisu rosbag2. Model składa się z dwóch części:
\begin{enumerate}[nosep]
  \item modelu czasu -- poprawienie znaczników czasu zegara systemowego
        na czas UTC (podrozdział~\ref{sec:model_czasu}),
  \item modelu odległości -- obliczenie odległości od stacji do samolotu
        (podrozdział~\ref{sec:model_odleglosc}).
\end{enumerate}

\subsection{Model czasu}
\label{sec:model_czasu}

Niech $t$ oznacza wskazanie zegara systemowego komputera, a $\tau$ --
odpowiadający mu rzeczywisty czas UTC. Zależność między nimi zapisano jako
\begin{equation}
  \tau = t + \delta,
  \label{eq:model_zegara}
\end{equation}
gdzie $\delta$ jest przesunięciem (błędem) zegara systemowego w sekundach.
Przyjęto, że w trakcie jednej sesji pomiarowej przesunięcie $\delta$ jest
stałe.

Węzeł obsługi GPS publikuje w temacie \texttt{/gps/time\_ref} pary
$(t_k, \tau_k)$, $k = 1, 2, \ldots, K$, gdzie $\tau_k$ jest czasem UTC
odczytanym z komunikatu \texttt{RMC}, a $t_k$ -- wskazaniem zegara
systemowego w chwili odczytu tego komunikatu. Każda para daje jeden pomiar
przesunięcia:
\begin{equation}
  \delta_k = \tau_k - t_k, \qquad k = 1, 2, \ldots, K.
  \label{eq:offset_pomiar}
\end{equation}
Pojedyncze pomiary różnią się nieco od siebie, dlatego jako wynik przyjmuje
się ich średnią arytmetyczną, a rozrzut pomiarów opisuje się odchyleniem
standardowym:
\begin{equation}
  \hat{\delta} = \frac{1}{K} \sum_{k=1}^{K} \delta_k,
  \qquad
  \sigma_{\delta} = \sqrt{\frac{1}{K-1} \sum_{k=1}^{K}
                    \bigl(\delta_k - \hat{\delta}\bigr)^2},
  \label{eq:offset_srednia}
\end{equation}
\begin{objasnienia}
  \item[$\hat{\delta}$] -- wyznaczone przesunięcie zegara w sekundach,
  \item[$\sigma_{\delta}$] -- odchylenie standardowe pomiarów przesunięcia w sekundach,
  \item[$K$] -- liczba pomiarów (przy jednym pomiarze na sekundę sesja trwająca \SI{10}{min} daje ok.~600 pomiarów).
\end{objasnienia}
Poprawiony czas UTC dowolnej zapisanej wiadomości o znaczniku $t$ wynosi
\begin{equation}
  \hat{\tau} = t + \hat{\delta}.
  \label{eq:korekta_czasu}
\end{equation}

Pomiar \eqref{eq:offset_pomiar} nie jest idealny, ponieważ znacznik $t_k$
jest nadawany dopiero po przesłaniu komunikatu łączem szeregowym i odczytaniu
go przez program. Przy prędkości \num{9600}~bit/s przesłanie jednego znaku
(10~bitów) trwa ok.\ \SI{1.04}{ms}, a komunikat NMEA ma do 82 znaków, więc
sama transmisja zajmuje do ok.\ \SI{85}{ms}. Program odczytuje port co
\SI{100}{ms}, co dodaje średnio kolejne \SI{50}{ms}. Wyznaczone przesunięcie
$\hat{\delta}$ jest więc obarczone błędem rzędu \SI{0.1}{s}, co -- jak
pokazano w podrozdziale~\ref{sec:model_przyklad} -- jest wartością
akceptowalną.

\subsection{Model odległości}
\label{sec:model_odleglosc}

\paragraph{Dane wejściowe.}
Z komunikatu SBS typu \texttt{MSG,3} odczytuje się szerokość $\varphi_i$
i długość $\lambda_i$ geograficzną samolotu (w stopniach) oraz jego wysokość
w stopach. Z tematu \texttt{/gps/fix} odczytuje się szerokość
$\varphi_{\mathrm{s}}$, długość $\lambda_{\mathrm{s}}$ i wysokość
$h_{\mathrm{s}}$ stacji (w metrach). Indeks $i$ oznacza samolot, a indeks
$\mathrm{s}$ -- stację.

\paragraph{Przeliczenie jednostek.}
W lotnictwie wysokość podaje się w stopach, a prędkość w węzłach. Na
jednostki układu SI przelicza się je wzorami
\begin{equation}
  h_i = \num{0.3048}\, h_{\mathrm{ft}}, \qquad
  v = \num{0.514444}\, v_{\mathrm{kt}},
  \label{eq:jednostki}
\end{equation}
\begin{objasnienia}
  \item[$h_{\mathrm{ft}}$] -- wysokość samolotu w stopach,
  \item[$h_i$] -- wysokość samolotu w metrach,
  \item[$v_{\mathrm{kt}}$] -- prędkość w węzłach,
  \item[$v$] -- prędkość w metrach na sekundę.
\end{objasnienia}

\paragraph{Odległość po powierzchni Ziemi.}
Ziemię przybliżono kulą o promieniu $R_{\mathrm{Z}} = \SI{6371}{km}$.
Odległość między stacją a punktem na powierzchni Ziemi leżącym dokładnie pod
samolotem oblicza się wzorem haversine:
\begin{equation}
  s = 2 R_{\mathrm{Z}} \arcsin \sqrt{
        \sin^2\frac{\Delta\varphi}{2}
        + \cos\varphi_{\mathrm{s}} \cos\varphi_i
          \sin^2\frac{\Delta\lambda}{2}},
  \label{eq:haversine}
\end{equation}
\begin{objasnienia}
  \item[$s$] -- odległość mierzona po powierzchni Ziemi w metrach,
  \item[$\Delta\varphi$] -- różnica szerokości geograficznych, $\Delta\varphi = \varphi_i - \varphi_{\mathrm{s}}$, w radianach,
  \item[$\Delta\lambda$] -- różnica długości geograficznych, $\Delta\lambda = \lambda_i - \lambda_{\mathrm{s}}$, w radianach,
  \item[$R_{\mathrm{Z}}$] -- średni promień Ziemi w metrach.
\end{objasnienia}
Kąty podawane w stopniach należy przed użyciem we wzorze
\eqref{eq:haversine} zamienić na radiany, mnożąc je przez $\pi/180$.

\paragraph{Odległość skośna.}
Samolot znajduje się na wysokości $h_i$, a stacja na wysokości
$h_{\mathrm{s}}$. Odległość w linii prostej od stacji do samolotu
(odległość skośną) przybliża się, traktując odległość $s$ i różnicę wysokości
jako przyprostokątne trójkąta prostokątnego:
\begin{equation}
  \rho = \sqrt{s^2 + \Delta h^2}, \qquad \Delta h = h_i - h_{\mathrm{s}},
  \label{eq:odleglosc_skosna}
\end{equation}
gdzie $\rho$ jest odległością skośną, a $\Delta h$ -- różnicą wysokości
samolotu i stacji, obie w metrach. Wzory \eqref{eq:jednostki}%
--\eqref{eq:odleglosc_skosna} są zaimplementowane w węźle wyświetlacza
stacji, który na ich podstawie pokazuje odległość do każdego samolotu.

\subsection{Uproszczenia modelu}
\label{sec:model_uproszczenia}

W modelu przyjęto następujące uproszczenia:
\begin{itemize}[nosep]
  \item Ziemia jest kulą, a nie elipsoidą -- powoduje to błąd odległości
        rzędu kilku dziesiątych procenta,
  \item we wzorze \eqref{eq:odleglosc_skosna} pominięto krzywiznę Ziemi;
        błąd rośnie wraz z odległością samolotu od stacji,
  \item wysokość barometryczna z ADS-B jest traktowana jak wysokość nad
        poziomem morza; przy ciśnieniu innym niż standardowe
        (\SI{1013.25}{hPa}) różnica może sięgać kilkudziesięciu metrów,
  \item przesunięcie zegara $\delta$ jest stałe w trakcie sesji (pominięto
        powolny dryf zegara),
  \item stacja jest nieruchoma w trakcie sesji,
  \item pominięto opóźnienie między odebraniem komunikatu ADS-B a nadaniem mu
        znacznika czasu (pojedyncze milisekundy).
\end{itemize}
Uproszczenia te są wystarczające do podglądu na żywo i do analizy zasięgu
stacji. Przy dalszym rozwoju systemu model odległości można zastąpić
dokładniejszym przeliczeniem współrzędnych do lokalnego układu kartezjańskiego
związanego ze stacją, które uwzględnia kształt Ziemi i pozwala dodatkowo
wyznaczyć kierunek obserwacji samolotu.

\subsection{Przykład obliczeniowy}
\label{sec:model_przyklad}

Przyjęto stację w punkcie $\varphi_{\mathrm{s}} = \SI{50.0}{\degree}$,
$\lambda_{\mathrm{s}} = \SI{20.0}{\degree}$, $h_{\mathrm{s}} = \SI{250}{m}$
oraz samolot w punkcie $\varphi_i = \SI{50.3169}{\degree}$,
$\lambda_i = \SI{20.49795}{\degree}$, lecący na wysokości
$h_{\mathrm{ft}} = \num{32800}$~stóp.

Z wzoru \eqref{eq:jednostki} wysokość samolotu wynosi
\begin{equation}
  h_i = \num{0.3048} \cdot \num{32800} \approx \SI{9997.4}{m},
  \qquad
  \Delta h = \num{9997.4} - 250 = \SI{9747.4}{m}.
\end{equation}
Różnice współrzędnych wynoszą $\Delta\varphi = \SI{0.3169}{\degree}$
i $\Delta\lambda = \SI{0.49795}{\degree}$. Wyrażenie pod pierwiastkiem we
wzorze \eqref{eq:haversine} ma wartość ok.\ \num{1.540e-5}, skąd
\begin{equation}
  s \approx \SI{50.00}{km}, \qquad
  \rho = \sqrt{\num{50000}^2 + \num{9747.4}^2} \approx \SI{50.94}{km}.
\end{equation}

Wpływ błędu czasu pokazuje następujące oszacowanie: samolot lecący
z prędkością $v = \SI{230}{m/s}$ w czasie $\Delta\tau$ pokonuje drogę
\begin{equation}
  \Delta x = v \, \Delta\tau .
  \label{eq:blad_czasu}
\end{equation}
Jeżeli zegar stacji myli się o \SI{10}{s} i nie zostanie poprawiony wzorem
\eqref{eq:korekta_czasu}, pozycja przypisana danej chwili będzie błędna
o ok.\ \SI{2.3}{km}. Po korekcji błąd czasu spada do ok.\ \SI{0.1}{s},
a błąd pozycji -- do ok.\ \SI{23}{m}. Uzasadnia to rejestrowanie czasu GPS
razem z danymi ADS-B.

\subsection{Weryfikacja modelu}
\label{sec:model_weryfikacja}

Poprawność modelu odległości można sprawdzić na prostych przypadkach:
dla samolotu położonego dokładnie nad stacją powinno zachodzić $s = 0$
i $\rho = \Delta h$, a dla punktów odległych o \SI{1}{\degree} szerokości
geograficznej $s \approx \SI{111.2}{km}$. Poprawność modelu czasu można
sprawdzić, porównując $\hat{\delta}$ wyznaczone w sesji z dostępem do
internetu (gdy zegar jest synchronizowany przez NTP i $\hat{\delta}$ powinno
być bliskie zera) z sesją bez dostępu do sieci. Szczegółową ocenę na
podstawie zarejestrowanych danych przedstawiono w~rozdziale~\ref{chap:badania}.
%  Wymaga w preambule: amsmath, siunitx, enumitem
% ============================================================

% Objaśnienia symboli pod wzorem ("gdzie:"): symbol w kolumnie od
% wcięcia akapitu, myślniki zawsze w jednej pionowej linii
\newenvironment{objasnienia}{%
  \noindent gdzie:
  \begin{description}[font=\normalfont, labelindent=\parindent,
    labelwidth=1.2cm, labelsep=0.2cm, leftmargin=!, align=left,
    itemsep=0pt, topsep=2pt]}{\end{description}}

\section{Model wyznaczania danych referencyjnych}
\label{sec:model}

Zarejestrowane przez stację dane mają posłużyć jako dane referencyjne, czyli
mają pozwolić stwierdzić, gdzie znajdował się dany statek powietrzny
w określonej chwili czasu. % referencja do podrozdziału.
W tym celu potrzebny jest model, który łączy trzy źródła informacji:
komunikaty ADS-B odbierane drogą radiową, pozycję stacji z odbiornika GPS oraz
zegar systemowy komputera, którym oznaczane są wszystkie zapisane wiadomości.
W niniejszym podrozdziale model ten przedstawiono najpierw w postaci słownej
(model opisowy), a następnie w postaci wzorów (model matematyczny). Na końcu
omówiono przyjęte uproszczenia i podano przykład obliczeniowy.

\subsection{Model opisowy}
\label{sec:model_opisowy}

Statek powietrzny wyposażony w transponder ADS-B kilka razy na sekundę
rozgłasza na częstotliwości \SI{1090}{MHz} komunikaty zawierające m.in.\ jego
unikalny adres ICAO, pozycję (szerokość i długość geograficzną) wyznaczoną
przez pokładowy odbiornik GNSS, wysokość barometryczną oraz prędkość
względem ziemi~\cite{sun2021}. Stacja odbiera część tych komunikatów --
niektóre giną z powodu zakłóceń, zbyt słabego sygnału lub przesłonięcia
anteny. Odebrane komunikaty są zapisywane w postaci tekstowej (format SBS)
razem ze znacznikiem czasu nadawanym przez komputer w chwili odbioru.

Komunikat ADS-B mówi, gdzie samolot jest, ale nie mówi, gdzie jest
względem stacji. Do tego potrzebna jest pozycja samej stacji, którą
podaje odbiornik GPS w komunikacie NMEA \texttt{GGA}~\cite{nmea0183}.
Znając obie pozycje, można obliczyć odległość od stacji do samolotu.

Drugim problemem jest czas. Komputer Raspberry Pi nie ma podtrzymywanego
bateryjnie zegara, a w terenie nie ma dostępu do internetu, więc jego zegar
może się mylić o nieznaną wartość. Ponieważ wszystkie zapisane wiadomości są
oznaczone właśnie tym zegarem, błąd zegara przenosi się na wszystkie dane.
Odbiornik GPS co sekundę przesyła jednak dokładny czas UTC (komunikat
\texttt{RMC}). Komputer zapisuje go razem z własnym wskazaniem zegara, dzięki
czemu po zakończeniu pomiarów można zmierzyć, o ile zegar się mylił,
i poprawić czas wszystkich wiadomości.

Granice modelowanego układu wyznacza sama stacja (antena, odbiornik RTL-SDR,
odbiornik GPS i komputer). Samoloty i satelity GPS należą do otoczenia,
a punktami styku są: łącze radiowe \SI{1090}{MHz}, łącze szeregowe
odbiornika GPS oraz plik zapisu rosbag2. Model składa się z dwóch części:
\begin{enumerate}[nosep]
  \item modelu czasu -- poprawienie znaczników czasu zegara systemowego
        na czas UTC (podrozdział~\ref{sec:model_czasu}),
  \item modelu odległości -- obliczenie odległości od stacji do samolotu
        (podrozdział~\ref{sec:model_odleglosc}).
\end{enumerate}

\subsection{Model czasu}
\label{sec:model_czasu}

Niech $t$ oznacza wskazanie zegara systemowego komputera, a $\tau$ --
odpowiadający mu rzeczywisty czas UTC. Zależność między nimi zapisano jako
\begin{equation}
  \tau = t + \delta,
  \label{eq:model_zegara}
\end{equation}
gdzie $\delta$ jest przesunięciem (błędem) zegara systemowego w sekundach.
Przyjęto, że w trakcie jednej sesji pomiarowej przesunięcie $\delta$ jest
stałe.

Węzeł obsługi GPS publikuje w temacie \texttt{/gps/time\_ref} pary
$(t_k, \tau_k)$, $k = 1, 2, \ldots, K$, gdzie $\tau_k$ jest czasem UTC
odczytanym z komunikatu \texttt{RMC}, a $t_k$ -- wskazaniem zegara
systemowego w chwili odczytu tego komunikatu. Każda para daje jeden pomiar
przesunięcia:
\begin{equation}
  \delta_k = \tau_k - t_k, \qquad k = 1, 2, \ldots, K.
  \label{eq:offset_pomiar}
\end{equation}
Pojedyncze pomiary różnią się nieco od siebie, dlatego jako wynik przyjmuje
się ich średnią arytmetyczną, a rozrzut pomiarów opisuje się odchyleniem
standardowym:
\begin{equation}
  \hat{\delta} = \frac{1}{K} \sum_{k=1}^{K} \delta_k,
  \qquad
  \sigma_{\delta} = \sqrt{\frac{1}{K-1} \sum_{k=1}^{K}
                    \bigl(\delta_k - \hat{\delta}\bigr)^2},
  \label{eq:offset_srednia}
\end{equation}
\begin{objasnienia}
  \item[$\hat{\delta}$] -- wyznaczone przesunięcie zegara w sekundach,
  \item[$\sigma_{\delta}$] -- odchylenie standardowe pomiarów przesunięcia w sekundach,
  \item[$K$] -- liczba pomiarów (przy jednym pomiarze na sekundę sesja trwająca \SI{10}{min} daje ok.~600 pomiarów).
\end{objasnienia}
Poprawiony czas UTC dowolnej zapisanej wiadomości o znaczniku $t$ wynosi
\begin{equation}
  \hat{\tau} = t + \hat{\delta}.
  \label{eq:korekta_czasu}
\end{equation}

Pomiar \eqref{eq:offset_pomiar} nie jest idealny, ponieważ znacznik $t_k$
jest nadawany dopiero po przesłaniu komunikatu łączem szeregowym i odczytaniu
go przez program. Przy prędkości \num{9600}~bit/s przesłanie jednego znaku
(10~bitów) trwa ok.\ \SI{1.04}{ms}, a komunikat NMEA ma do 82 znaków, więc
sama transmisja zajmuje do ok.\ \SI{85}{ms}. Program odczytuje port co
\SI{100}{ms}, co dodaje średnio kolejne \SI{50}{ms}. Wyznaczone przesunięcie
$\hat{\delta}$ jest więc obarczone błędem rzędu \SI{0.1}{s}, co -- jak
pokazano w podrozdziale~\ref{sec:model_przyklad} -- jest wartością
akceptowalną.

\subsection{Model odległości}
\label{sec:model_odleglosc}

\paragraph{Dane wejściowe.}
Z komunikatu SBS typu \texttt{MSG,3} odczytuje się szerokość $\varphi_i$
i długość $\lambda_i$ geograficzną samolotu (w stopniach) oraz jego wysokość
w stopach. Z tematu \texttt{/gps/fix} odczytuje się szerokość
$\varphi_{\mathrm{s}}$, długość $\lambda_{\mathrm{s}}$ i wysokość
$h_{\mathrm{s}}$ stacji (w metrach). Indeks $i$ oznacza samolot, a indeks
$\mathrm{s}$ -- stację.

\paragraph{Przeliczenie jednostek.}
W lotnictwie wysokość podaje się w stopach, a prędkość w węzłach. Na
jednostki układu SI przelicza się je wzorami
\begin{equation}
  h_i = \num{0.3048}\, h_{\mathrm{ft}}, \qquad
  v = \num{0.514444}\, v_{\mathrm{kt}},
  \label{eq:jednostki}
\end{equation}
\begin{objasnienia}
  \item[$h_{\mathrm{ft}}$] -- wysokość samolotu w stopach,
  \item[$h_i$] -- wysokość samolotu w metrach,
  \item[$v_{\mathrm{kt}}$] -- prędkość w węzłach,
  \item[$v$] -- prędkość w metrach na sekundę.
\end{objasnienia}

\paragraph{Odległość po powierzchni Ziemi.}
Ziemię przybliżono kulą o promieniu $R_{\mathrm{Z}} = \SI{6371}{km}$.
Odległość między stacją a punktem na powierzchni Ziemi leżącym dokładnie pod
samolotem oblicza się wzorem haversine:
\begin{equation}
  s = 2 R_{\mathrm{Z}} \arcsin \sqrt{
        \sin^2\frac{\Delta\varphi}{2}
        + \cos\varphi_{\mathrm{s}} \cos\varphi_i
          \sin^2\frac{\Delta\lambda}{2}},
  \label{eq:haversine}
\end{equation}
\begin{objasnienia}
  \item[$s$] -- odległość mierzona po powierzchni Ziemi w metrach,
  \item[$\Delta\varphi$] -- różnica szerokości geograficznych, $\Delta\varphi = \varphi_i - \varphi_{\mathrm{s}}$, w radianach,
  \item[$\Delta\lambda$] -- różnica długości geograficznych, $\Delta\lambda = \lambda_i - \lambda_{\mathrm{s}}$, w radianach,
  \item[$R_{\mathrm{Z}}$] -- średni promień Ziemi w metrach.
\end{objasnienia}
Kąty podawane w stopniach należy przed użyciem we wzorze
\eqref{eq:haversine} zamienić na radiany, mnożąc je przez $\pi/180$.

\paragraph{Odległość skośna.}
Samolot znajduje się na wysokości $h_i$, a stacja na wysokości
$h_{\mathrm{s}}$. Odległość w linii prostej od stacji do samolotu
(odległość skośną) przybliża się, traktując odległość $s$ i różnicę wysokości
jako przyprostokątne trójkąta prostokątnego:
\begin{equation}
  \rho = \sqrt{s^2 + \Delta h^2}, \qquad \Delta h = h_i - h_{\mathrm{s}},
  \label{eq:odleglosc_skosna}
\end{equation}
gdzie $\rho$ jest odległością skośną, a $\Delta h$ -- różnicą wysokości
samolotu i stacji, obie w metrach. Wzory \eqref{eq:jednostki}%
--\eqref{eq:odleglosc_skosna} są zaimplementowane w węźle wyświetlacza
stacji, który na ich podstawie pokazuje odległość do każdego samolotu.

\subsection{Uproszczenia modelu}
\label{sec:model_uproszczenia}

W modelu przyjęto następujące uproszczenia:
\begin{itemize}[nosep]
  \item Ziemia jest kulą, a nie elipsoidą -- powoduje to błąd odległości
        rzędu kilku dziesiątych procenta,
  \item we wzorze \eqref{eq:odleglosc_skosna} pominięto krzywiznę Ziemi;
        błąd rośnie wraz z odległością samolotu od stacji,
  \item wysokość barometryczna z ADS-B jest traktowana jak wysokość nad
        poziomem morza; przy ciśnieniu innym niż standardowe
        (\SI{1013.25}{hPa}) różnica może sięgać kilkudziesięciu metrów,
  \item przesunięcie zegara $\delta$ jest stałe w trakcie sesji (pominięto
        powolny dryf zegara),
  \item stacja jest nieruchoma w trakcie sesji,
  \item pominięto opóźnienie między odebraniem komunikatu ADS-B a nadaniem mu
        znacznika czasu (pojedyncze milisekundy).
\end{itemize}
Uproszczenia te są wystarczające do podglądu na żywo i do analizy zasięgu
stacji. Przy dalszym rozwoju systemu model odległości można zastąpić
dokładniejszym przeliczeniem współrzędnych do lokalnego układu kartezjańskiego
związanego ze stacją, które uwzględnia kształt Ziemi i pozwala dodatkowo
wyznaczyć kierunek obserwacji samolotu.

\subsection{Przykład obliczeniowy}
\label{sec:model_przyklad}

Przyjęto stację w punkcie $\varphi_{\mathrm{s}} = \SI{50.0}{\degree}$,
$\lambda_{\mathrm{s}} = \SI{20.0}{\degree}$, $h_{\mathrm{s}} = \SI{250}{m}$
oraz samolot w punkcie $\varphi_i = \SI{50.3169}{\degree}$,
$\lambda_i = \SI{20.49795}{\degree}$, lecący na wysokości
$h_{\mathrm{ft}} = \num{32800}$~stóp.

Z wzoru \eqref{eq:jednostki} wysokość samolotu wynosi
\begin{equation}
  h_i = \num{0.3048} \cdot \num{32800} \approx \SI{9997.4}{m},
  \qquad
  \Delta h = \num{9997.4} - 250 = \SI{9747.4}{m}.
\end{equation}
Różnice współrzędnych wynoszą $\Delta\varphi = \SI{0.3169}{\degree}$
i $\Delta\lambda = \SI{0.49795}{\degree}$. Wyrażenie pod pierwiastkiem we
wzorze \eqref{eq:haversine} ma wartość ok.\ \num{1.540e-5}, skąd
\begin{equation}
  s \approx \SI{50.00}{km}, \qquad
  \rho = \sqrt{\num{50000}^2 + \num{9747.4}^2} \approx \SI{50.94}{km}.
\end{equation}

Wpływ błędu czasu pokazuje następujące oszacowanie: samolot lecący
z prędkością $v = \SI{230}{m/s}$ w czasie $\Delta\tau$ pokonuje drogę
\begin{equation}
  \Delta x = v \, \Delta\tau .
  \label{eq:blad_czasu}
\end{equation}
Jeżeli zegar stacji myli się o \SI{10}{s} i nie zostanie poprawiony wzorem
\eqref{eq:korekta_czasu}, pozycja przypisana danej chwili będzie błędna
o ok.\ \SI{2.3}{km}. Po korekcji błąd czasu spada do ok.\ \SI{0.1}{s},
a błąd pozycji -- do ok.\ \SI{23}{m}. Uzasadnia to rejestrowanie czasu GPS
razem z danymi ADS-B.

\subsection{Weryfikacja modelu}
\label{sec:model_weryfikacja}

Poprawność modelu odległości można sprawdzić na prostych przypadkach:
dla samolotu położonego dokładnie nad stacją powinno zachodzić $s = 0$
i $\rho = \Delta h$, a dla punktów odległych o \SI{1}{\degree} szerokości
geograficznej $s \approx \SI{111.2}{km}$. Poprawność modelu czasu można
sprawdzić, porównując $\hat{\delta}$ wyznaczone w sesji z dostępem do
internetu (gdy zegar jest synchronizowany przez NTP i $\hat{\delta}$ powinno
być bliskie zera) z sesją bez dostępu do sieci. Szczegółową ocenę na
podstawie zarejestrowanych danych przedstawiono w~rozdziale~\ref{chap:badania}.         % 3. Wymagania, architektura, tematy, model (35_Model)
\chapter{Warstwa sprzętowa stacji bazowej}
\label{chap:sprzet}


\section{Dobór podzespołów}
\label{sec:dobor}
% ANALIZA porównawcza na podstawie danych katalogowych

Dobór podzespołów przeprowadzono na podstawie analizy porównawczej parametrów
katalogowych dostępnych rozwiązań, przyjmując kryteria wynikające z~wymagań
przedstawionych w~podrozdziale~\ref{sec:wymagania}. Jako dodatkowe kryterium
uwzględniono dostępność podzespołów, co pozwoliło skrócić czas budowy
prototypu i~ograniczyć jego koszt. Poprawność dokonanego wyboru zweryfikowano
pomiarami przedstawionymi w~rozdziale~\ref{chap:badania}.

\begin{table}[H]
  \centering
  \caption{Porównanie platform obliczeniowych}
  \label{tab:platformy}
  \small
  \begin{tabular}{l c c c c c}
    \toprule
    Platforma & Pobór mocy [W] & GPIO/UART/SPI & ROS~2 Jazzy & GPU & Koszt \\
    \midrule
    Raspberry Pi 4B       & -- & -- & -- & -- & -- \\
    Raspberry Pi 5        & -- & -- & -- & -- & -- \\
    NVIDIA Jetson Orin Nano & -- & -- & -- & -- & -- \\
    Mini-PC x86           & -- & -- & -- & -- & -- \\
    \bottomrule
  \end{tabular}
\end{table}

\begin{table}[H]
  \centering
  \caption{Porównanie odbiorników SDR}
  \label{tab:sdr}
  \small
  \begin{tabular}{l c c c c c}
    \toprule
    Odbiornik & Zakres [MHz] & ADC [bit] & Filtr 1090~MHz & dump1090 & Koszt \\
    \midrule
    RTL-SDR Blog V3             & -- & -- & -- & -- & -- \\
    FlightAware Pro Stick Plus  & -- & -- & -- & -- & -- \\
    Airspy Mini                 & -- & -- & -- & -- & -- \\
    HackRF One                  & -- & -- & -- & -- & -- \\
    \bottomrule
  \end{tabular}
\end{table}

Wybrane podzespoły zestawiono w~tabeli~\ref{tab:podzespoly}.

\begin{table}[H]
  \centering
  \caption{Zestawienie podzespołów stacji bazowej}
  \label{tab:podzespoly}
  \small
  \begin{tabular}{l l >{\raggedright\arraybackslash}p{5.5cm}}
    \toprule
    Podzespół & Model & Uwagi \\
    \midrule
    Komputer        & Raspberry Pi 4B          & system uruchamiany z~pamięci USB \SI{64}{GB} \\
    Odbiornik SDR   & RTL-SDR Blog V3          & złącze USB \\
    Antena          & dipol z~zestawu RTL-SDR Blog & ramiona ok.\ \SI{6.5}{cm}, montaż pionowy \\
    Odbiornik GPS   & moduł z~wyjściem NMEA    & UART, \SI{9600}{bit/s}, komunikaty GGA i~RMC co \SI{1}{s} \\
    Wyświetlacz     & TFT ILI9341, $240\times320$ & przekątna 2,4~cala, magistrala SPI \\
    Zasilanie       & powerbank LYNXER RPP-272 & wyjście USB \SI{5}{V} \\
    \bottomrule
  \end{tabular}
\end{table}

% ---------------------------------------------------------------------
\section{Tor odbioru radiowego}
\label{sec:tor_radiowy}
% Antena dipolowa ćwierćfalowa (odwołanie do wzoru
% \eqref{eq:lambda}), polaryzacja pionowa, RTL-SDR, zakłócenia od USB,
% umiejscowienie anteny.

% ---------------------------------------------------------------------
\section{Odbiornik GPS}
\label{sec:odbiornik_gps}
%Połączenie UART, ustawienia, komunikaty, antena z widokiem
% na niebo, utrata fixa w pomieszczeniu.

% ---------------------------------------------------------------------
\section{Elementy obsługi i~schemat połączeń}
\label{sec:polaczenia}
% Wyświetlacz, dioda, przycisk; schemat (rysunek);
% konflikt GPIO3 z I2C (dlaczego I2C wyłączone).

Połączenia elementów z~wyprowadzeniami złącza GPIO komputera Raspberry Pi
zestawiono w~tabeli~\ref{tab:gpio}.

\begin{table}[H]
  \centering
  \caption{Połączenia ze złączem GPIO komputera Raspberry Pi}
  \label{tab:gpio}
  \small
  \begin{tabular}{>{\raggedright}p{3.6cm} >{\raggedright}p{5.2cm} >{\raggedright\arraybackslash}p{4.8cm}}
    \toprule
    Element & Wyprowadzenia (numer fizyczny) & Uwagi \\
    \midrule
    Dioda stanu          & 11 (GPIO17), 34 (GND) & rezystor szeregowy \SIrange{220}{330}{\ohm} \\
    Przycisk zasilania   & 5 (GPIO3), 6 (GND)    & wybudzanie ze stanu zatrzymania \\
    Wyświetlacz (SPI0)   & 19, 21, 23, 24; DC: 18 (GPIO24); RST: 22 (GPIO25) & sterownik ILI9341 \\
    Odbiornik GPS (UART) & 8, 10                 & \texttt{/dev/ttyAMA0} \\
    \bottomrule
  \end{tabular}
\end{table}

% ---------------------------------------------------------------------
\section{Zasilanie i~obudowa}
\label{sec:zasilanie}
% Powerbank, obudowa drukowana 3D, chłodzenie, okienko na
% wskaźnik powerbanka, dostęp do ładowania.               % 4. Warstwa sprzętowa stacji bazowej
% =====================================================================
%  5. WARSTWA PROGRAMOWA STACJI BAZOWEJ  (dawne 40_Implementation.tex)
%  Objętość: ok. 13 stron - największy rozdział (własna praca).
% =====================================================================
\chapter{Warstwa programowa stacji bazowej}
\label{chap:oprogramowanie}

% 2-3 zdania: co jest w tym rozdziale.

% ---------------------------------------------------------------------
\section{Środowisko uruchomieniowe}
\label{sec:srodowisko}
% (ok. 2 str.) Ubuntu, Docker (obraz, kontener, compose), usługi systemd,
% dump1090 na hoście (porty 30002 i 30003).

% ---------------------------------------------------------------------
\section{Węzły akwizycji danych}
\label{sec:akwizycja}
% (ok. 2 str.) sdr_node (dwa połączenia TCP, buforowanie linii,
% automatyczne ponowne łączenie), gps_node, power_monitor_node.
% ======================================================================
% rys_schemat_pozycja.tex – schemat blokowy wyznaczania pozycji statku
% powietrznego w module adsb_decoder.py (metoda _update_position)
% Wstawić w rozdziale 4 (Implementacja) poleceniem: \input{tex/rys_schemat_pozycja}
% Wymaga w preambule: tikz + biblioteki positioning, arrows.meta, calc,
%                     shapes.geometric (diamond, trapezium); siunitx; kolor akcent.
% Odwołuje się do tabeli tab:tablica_pozycja (plik tab_tablica_pozycja.tex).
% Treść zgodna z kodem w repozytorium (commit 33bb6ff).
% Kod TikZ przygotowano z wykorzystaniem Claude (Anthropic) – zob. deklarację AI.
% ======================================================================
\begin{figure}[!htb]
\centering
\begin{tikzpicture}[
    font=\scriptsize,
    >={Stealth[length=2mm]},
    term/.style={draw, thick, rounded corners=3.5mm, fill=akcent!25,
                 minimum height=7mm, minimum width=22mm, align=center},
    we/.style={draw, thick, trapezium, trapezium left angle=70,
               trapezium right angle=110, fill=black!8, align=center,
               minimum height=9mm, inner xsep=4mm, text width=40mm},
    op/.style={draw, thick, fill=akcent!10, align=center,
               minimum height=8mm, text width=54mm, inner sep=3pt},
    odrz/.style={draw, thick, fill=akcent!10, align=center,
                 minimum height=8mm, text width=31mm, inner sep=3pt},
    dec/.style={draw, thick, diamond, aspect=2.6, fill=white, align=center,
                inner sep=0pt, text width=42mm},
    strz/.style={->, thick},
    tn/.style={font=\scriptsize\bfseries, inner sep=1.5pt},
]
\hyphenpenalty=10000 \exhyphenpenalty=10000
% ---------------- główna ścieżka ----------------
\node[term] (start) at (0, 0) {START};
\node[we]   (wej)   at (0,-1.35) {Ramka DF17 z~pozycją, czas $t$};
\node[dec]  (d0)    at (0,-3.2) {Znana pozycja stacji?};
\node[op]   (opl)   at (0,-5.3)
     {$L\leftarrow$ pozycja lokalna z~jednej ramki (odniesienie: stacja)};
\node[op]   (opg)   at (0,-6.8)
     {$G\leftarrow$ pozycja globalna z~pary ramek (PipeDecoder) lub brak};
\node[op]   (tab)   at (0,-8.3)
     {Ustal $P$ i~$Q$ według tablicy decyzyjnej (tabela~\ref{tab:tablica_pozycja})};
\node[dec]  (d1)    at (0,-10.3) {Ustalono pozycję $P$?};
\node[dec]  (d2)    at (0,-13.1)
     {Stacja znana i~$P$ dalej niż \SI{450}{km} od stacji?};
\node[dec]  (d3)    at (0,-16.5)
     {Skok od poprzedniej pozycji $>\SI{600}{m/s}\cdot\Delta t+\SI{1}{km}$
      i~$Q\neq$ VERIFIED?};
\node[op]   (zap)   at (0,-19.0) {Zapisz $P$, $Q$ oraz czas pozycji $t$};
\node[term] (konz)  at (0,-20.3) {KONIEC};

% ---------------- gałąź „brak pozycji stacji” ----------------
\node[op, text width=26mm] (lbrak) at (6.2,-3.2) {$L\leftarrow$ brak};

% ---------------- odrzucenia po prawej ----------------
\node[odrz] (r1) at (6.2,-10.3) {Odrzuć: brak pozycji lub niezgodne $L$ i~$G$};
\node[term] (k1) at (9.6,-10.3) {KONIEC};
\node[odrz] (r2) at (6.2,-13.1) {Odrzuć: poza zasięgiem};
\node[term] (k2) at (9.6,-13.1) {KONIEC};
\node[odrz] (r3) at (6.2,-16.5) {Odrzuć: nierealny skok pozycji};
\node[term] (k3) at (9.6,-16.5) {KONIEC};

% ---------------- strzałki: główna ścieżka ----------------
\draw[strz] (start) -- (wej);
\draw[strz] (wej) -- (d0);
\draw[strz] (d0) -- node[tn, right=1mm] {TAK} (opl);
\draw[strz] (d0) -- node[tn, above, pos=0.35] {NIE} (lbrak);
\draw[strz] (opl) -- (opg);
\draw[strz] (lbrak.south) |- (opg.east);
\draw[strz] (opg) -- (tab);
\draw[strz] (tab) -- (d1);
\draw[strz] (d1) -- node[tn, right=1mm] {TAK} (d2);
\draw[strz] (d2) -- node[tn, right=1mm] {NIE} (d3);
\draw[strz] (d3) -- node[tn, right=1mm] {NIE} (zap);
\draw[strz] (zap) -- (konz);

% ---------------- strzałki: odrzucenia ----------------
\draw[strz] (d1) -- node[tn, above, pos=0.3] {NIE} (r1);
\draw[strz] (r1) -- (k1);
\draw[strz] (d2) -- node[tn, above, pos=0.3] {TAK} (r2);
\draw[strz] (r2) -- (k2);
\draw[strz] (d3) -- node[tn, above, pos=0.3] {TAK} (r3);
\draw[strz] (r3) -- (k3);
\end{tikzpicture}
\caption{Schemat blokowy wyznaczania pozycji statku powietrznego w~module
\texttt{adsb\_decoder.py}. Symbole: owal -- początek i~koniec algorytmu,
równoległobok -- dane wejściowe, prostokąt -- operacja, romb -- decyzja
z~wyjściami TAK i~NIE. Oznaczenia: $L$ -- pozycja wyznaczona metodą lokalną,
$G$ -- metodą globalną, $P$ -- pozycja wynikowa, $Q$ -- jakość pozycji
(LOCAL, GLOBAL lub VERIFIED), $\Delta t$ -- czas od poprzedniej pozycji tego
statku. Schemat obejmuje etap następujący po odfiltrowaniu ramek, które nie
są ramkami DF17 lub mają błędną sumę CRC}
\label{rys:schemat_pozycja}
\end{figure}




%\begin{lstlisting}[language=Python, caption={Dzielenie strumienia TCP na linie SBS}, %label={lst:linie}]
%self.buffer += data
%while b'\n' in self.buffer:
%    line, self.buffer = self.buffer.split(b'\n', 1)
%    # publikujemy tylko komunikaty zaczynające się od MSG
%\end{lstlisting}

%Na listingu~\ref{lst:linie} pokazano\ldots
% TODO: zaktualizuj listing do bieżącej wersji sdr_node.py (klasa TcpLineClient).

% ---------------------------------------------------------------------
\section{Dekodowanie ADS-B w~czasie rzeczywistym}
\label{sec:dekodowanie}
% (ok. 3,5 str.) Moduł adsb_decoder, węzeł decoder_node, dwie metody
% pozycji (lokalna z pozycją stacji, globalna z pary ramek), statusy
% LOCAL/GLOBAL/VERIFIED, kontrole wiarygodności, wyznaczanie odległości,
% azymutu i elewacji, publikacja AircraftState.

% ---------------------------------------------------------------------
\section{Rejestracja danych}
\label{sec:rejestracja}
% (ok. 2 str.) rosbag2 + MCAP, sesje nazywane czasem UTC, pliki co
% 10 min, odzyskiwanie po utracie zasilania (reindex).

% ---------------------------------------------------------------------
\section{Interfejs użytkownika}
\label{sec:interfejs}
% (ok. 1,5 str.) display_node (układ ekranu, kolory jakości pozycji,
% wskaźnik REC), skrypt przycisku zasilania.

% ---------------------------------------------------------------------
\section{Uruchamianie i~zamykanie systemu}
\label{sec:start_stop}
% (ok. 2 str.) start_nodes.sh, sygnały SIGTERM/SIGINT, trap, grupy
% procesów, SIGTTIN; napotkane problemy i ich rozwiązania.               % 5. Warstwa programowa stacji bazowej
% =====================================================================
%  6. BADANIA STACJI BAZOWEJ  (wersja 2026-10-29)
%  Objętość: ok. 10 stron. Podział jak w rozdz. 4 i 5: sprzęt / program.
% =====================================================================
\chapter{Badania stacji bazowej}
\label{chap:badania}

% 2-3 zdania: co jest w tym rozdziale.

% ---------------------------------------------------------------------
\section{Metodyka i~scenariusze badań}
\label{sec:metodyka}
% (ok. 1,5 str.) Cele badań, scenariusze (sesje: data, miejsce, czas
% trwania), sposób analizy danych z rosbag2, odwołania do wzorów modelu.

% ---------------------------------------------------------------------
\section{Badania warstwy sprzętowej}
\label{sec:badania_sprzet}

\subsection{Zasięg odbioru}
\label{sec:zasieg}
% Maksymalna odległość, mapa/rozkład odległości, wpływ anteny.

% ---------------------------------------------------------------------
% Podrozdzial: czas pracy na zasilaniu bateryjnym (test T6.1)
% Wkleic ZAMIAST trzech linii:
%   \subsection{Czas pracy na zasilaniu bateryjnym}
%   \label{sec:czas_pracy}
%   % Test od 100 %, porównanie z pojemnością powerbanka.
% Miejsca do uzupelnienia: \uzup{...} (czerwony tekst). Definicje makra
% \uzup mozna przeniesc do 00_preamble.tex i usunac na koniec pisania.
% ---------------------------------------------------------------------
\providecommand{\uzup}[1]{\textcolor{red}{[#1]}}

\subsection{Czas pracy na zasilaniu bateryjnym}
\label{sec:czas_pracy}

Stacja bazowa ma pracować w~terenie bez dostępu do sieci energetycznej, dlatego jednym z~podstawowych parametrów jest czas pracy na jednym naładowaniu powerbanku. Celem testu było wyznaczenie tego czasu przy pełnym obciążeniu systemu i~porównanie go z~pojemnością powerbanku. Wymagany czas pracy wynosi co najmniej \SI{8}{h}.

Badano kompletną stację bazową zasilaną z~powerbanku LYNXER RPP-272 przez złącze USB-C. Podczas testu działały wszystkie węzły ROS~2 oraz rejestracja sesji rosbag2. Pomiar wykonano w~dniu \uzup{data} w~miejscu \uzup{opis miejsca}, przy temperaturze otoczenia \uzup{T\_ot}~\si{\degreeCelsius}. Stacja była \uzup{zamknięta w~obudowie z~wentylatorem / otwarta na stanowisku}. Zegar systemowy był synchronizowany przez NTP, co zapewniło poprawne znaczniki czasu w~dzienniku systemowym.

Test przeprowadzono według następujących kroków:
\begin{enumerate}
  \item naładowanie powerbanku do \SI{100}{\percent} według jego wskaźnika i~odłączenie ładowarki,
  \item podłączenie stacji do powerbanku, co powoduje uruchomienie systemu i~automatyczny start nowej sesji pomiarowej,
  \item notowanie w~odstępach około \SI{1}{h} wskazania powerbanku, temperatury procesora i~liczby śledzonych statków powietrznych z~wyświetlacza stacji,
  \item pozostawienie stacji bez ingerencji aż do samoczynnego odcięcia zasilania przez powerbank,
  \item odczytanie chwili uruchomienia i~chwili ostatniego wpisu z~dziennika systemowego.
\end{enumerate}

Czas pracy na baterii obliczono jako
\begin{equation}
  T_{\mathrm{bat}} = t_{\mathrm{k}} - t_{\mathrm{p}},
  \label{eq:t_bat}
\end{equation}
gdzie $t_{\mathrm{p}}$ oznacza chwilę uruchomienia systemu, a~$t_{\mathrm{k}}$ chwilę ostatniego wpisu w~dzienniku przed odcięciem zasilania. Aby porównać wynik z~pojemnością powerbanku, wyznaczono górne oszacowanie średniej mocy pobieranej przez stację:
\begin{equation}
  \bar{P} \leq \frac{E_{\mathrm{nom}}}{T_{\mathrm{bat}}}, \qquad E_{\mathrm{nom}} = Q_{\mathrm{nom}} \, U_{\mathrm{ogn}},
  \label{eq:p_sr}
\end{equation}
gdzie $E_{\mathrm{nom}}$ oznacza energię znamionową powerbanku, $Q_{\mathrm{nom}}$ jego deklarowaną pojemność, a~$U_{\mathrm{ogn}}$ napięcie znamionowe ogniw. Nierówność wynika ze strat w~przetwornicy powerbanku i~z~tego, że rzeczywista pojemność może być mniejsza od deklarowanej.

Wyniki pomiaru zestawiono w~tabeli~\ref{tab:bateria_wyniki}, a~obserwacje prowadzone w~trakcie testu w~tabeli~\ref{tab:bateria_obserwacje}.

\begin{table}[!htb]
  \centering
  \caption{Wyniki testu czasu pracy stacji bazowej na zasilaniu bateryjnym}
  \label{tab:bateria_wyniki}
  \begin{tabular}{lcc}
    \toprule
    Wielkość & Wartość & Wymaganie \\
    \midrule
    Uruchomienie systemu $t_{\mathrm{p}}$ (UTC) & \uzup{hh:mm:ss} & -- \\
    Ostatni wpis w~dzienniku $t_{\mathrm{k}}$ (UTC) & \uzup{hh:mm:ss} & -- \\
    Czas pracy $T_{\mathrm{bat}}$ [h] & \uzup{np. 11,2} & $\geq 8$ \\
    Energia znamionowa $E_{\mathrm{nom}}$ [Wh] & \uzup{z~etykiety} & -- \\
    Górne oszacowanie mocy $\bar{P}$ [W] & \uzup{wynik} & -- \\
    Zdarzenia niskiego napięcia [-] & \uzup{0} & -- \\
    \bottomrule
  \end{tabular}
\end{table}

\begin{table}[!htb]
  \centering
  \caption{Obserwacje prowadzone w~trakcie testu czasu pracy na baterii: $t$~-- czas od uruchomienia systemu, $B$~-- wskazanie powerbanku, $\vartheta_{\mathrm{CPU}}$~-- temperatura procesora, $N_{\mathrm{s}}$~-- liczba śledzonych statków powietrznych (odczyt z~wyświetlacza stacji)}
  \label{tab:bateria_obserwacje}
  \begin{tabular}{cccc}
    \toprule
    $t$ [h] & $B$ [\%] & $\vartheta_{\mathrm{CPU}}$ [\si{\degreeCelsius}] & $N_{\mathrm{s}}$ [-] \\
    \midrule
    0 & 100 & \uzup{} & \uzup{} \\
    1 & \uzup{} & \uzup{} & \uzup{} \\
    2 & \uzup{} & \uzup{} & \uzup{} \\
    3 & \uzup{} & \uzup{} & \uzup{} \\
    4 & \uzup{} & \uzup{} & \uzup{} \\
    5 & \uzup{} & \uzup{} & \uzup{} \\
    6 & \uzup{} & \uzup{} & \uzup{} \\
    7 & \uzup{} & \uzup{} & \uzup{} \\
    8 & \uzup{} & \uzup{} & \uzup{} \\
    \bottomrule
  \end{tabular}
\end{table}

Stacja pracowała na jednym naładowaniu powerbanku przez \uzup{T\_bat}~\si{h}, co \uzup{spełnia / nie spełnia} wymaganie minimalnego czasu pracy wynoszące \SI{8}{h}. Wynik jest \uzup{zgodny / niezgodny} z~wcześniejszym szacunkiem około \SIrange{11}{12}{h}, wyznaczonym przez ekstrapolację testu rozpoczętego przy \SI{35}{\percent} naładowania. Z~zależności~\eqref{eq:p_sr} wynika, że średnia moc pobierana przez stację nie przekraczała \uzup{P}~\si{W}. Wskazanie powerbanku malało \uzup{równomiernie / nierównomiernie}, a~ostatnie \uzup{X}~\si{\percent} odpowiadało \uzup{Y}~\si{min} pracy, dlatego wskaźnik ten nadaje się tylko do zgrubnej oceny pozostałego czasu.

Przed odcięciem zasilania zarejestrowano \uzup{0} zdarzeń niskiego napięcia. Potwierdza to, że użyty powerbank utrzymuje napięcie \SI{5}{V} do końca pracy i~wyłącza się bez ostrzeżenia, więc automatyczne zamykanie systemu na podstawie pomiaru napięcia nie byłoby skuteczne. Skutki takiego odcięcia dla zarejestrowanej sesji omówiono w~podrozdziale~\ref{sec:utrata_zasilania}, a~temperaturę procesora w~podrozdziale~\ref{sec:termika}.

\subsection{Warunki termiczne}
\label{sec:termika}
% Temperatura CPU na stole i w obudowie (temat /system/cpu_temp).

\subsection{Odporność na nagłą utratę zasilania}
\label{sec:utrata_zasilania}
% Test 2026-10-19/20: utracony ostatni plik, odzyskanie metadanych.

% ---------------------------------------------------------------------
\section{Badania warstwy programowej}
\label{sec:badania_program}

\subsection{Weryfikacja dekodera na danych syntetycznych}
\label{sec:dane_syntetyczne}
% Symulowany przelot, błąd pozycji, odrzucanie błędnych ramek,
% wykrycie błędnej pozycji stacji.

\subsection{Jakość wyznaczanych pozycji}
\label{sec:jakosc_pozycji}
% Czas do pierwszej pozycji (lokalna vs globalna), odsetek VERIFIED,
% odrzucone pozycje, porównanie z dump1090 (SBS).

\subsection{Synchronizacja czasu i~dostępność GNSS}
\label{sec:czas_gnss}
% Wyznaczone przesunięcie zegara (model z rozdz. 3), przerwy w fixie.

\subsection{Obciążenie procesora i~ilość rejestrowanych danych}
\label{sec:obciazenie}
% Obciążenie CPU przez węzły, przyrost danych w MB/h.   % 6. Badania stacji bazowej
\section*{Podsumowanie} \addcontentsline{toc}{section}{Podsumowanie}
osiagniecie celu
dalszy rozwoj             % 7. Podsumowanie i wnioski

% -------------------------- LITERATURA ------------------------------
\bibliographystyle{unsrt}        % numeracja w kolejności cytowania
\bibliography{bibliography}      % plik bibliography.bib

% --------------------------- ZAŁĄCZNIKI -----------------------------
\appendix
\chapter{Instrukcja uruchomienia i~obsługi stacji bazowej}
\label{app:instrukcja}
% Włączenie, odczyt wyświetlacza, wyłączenie przyciskiem,
% kopiowanie sesji na komputer, odzyskiwanie danych po utracie zasilania.

\chapter{Kod źródłowy}
\label{app:kod}

Pełny kod źródłowy dostępny jest w~repozytorium:
\url{https://github.com/1maks2/adsb-ros2-station}.

% Krótki opis najważniejszych plików (węzły, start_nodes.sh, Dockerfile).

\end{document}